% !TEX program = xelatex
% =============================================================================
% ЗАВРШНИ (ДИПЛОМСКИ) РАД — ћирилица
% Форматирање према: „Прилог 2а — Шаблон за обликовање завршног рада“ (ФОН)
% Мере су преузете из самог .docx шаблона (стилови и распоред страна):
%   - А4, маргине: лева 3,1 cm; горња, доња и десна 2,5 cm
%     (насловна страна, посвета, апстракт и садржај: 2,54 cm, доња 1,5 cm)
%   - текст: Times New Roman 12, обострано, проред 1,5 (20,7 pt), без растављања
%   - размак пасуса: 12 pt пре и 6 pt после (између пасуса 12 pt)
%   - наслови: Tahoma 14; 1. ниво ВЕЛИКА СЛОВА затамњено, 2. ниво ВЕЛИКА СЛОВА,
%     3. ниво мала слова; поглавља се не почињу на новој страни
%   - набрајање цртицом „-“, нумерисање „1.“
%   - табеле са пуним ивицама, „Табела 1: …“ изнад; „Слика 1. …“ испод (ознака
%     затамњена)
%   - нумерација страна од прве стране увода, десно у подножју
%
% ВАЖНО: компајлирати ИСКЉУЧИВО са XeLaTeX (не pdflatex!):
%   xelatex "Diplomski rad" ; biber "Diplomski rad" ; xelatex ; xelatex
% Потребни фонтови: Times New Roman и Tahoma (Tahoma је копирана са Windows
% партиције у ~/.local/share/fonts/Tahoma).
% =============================================================================

\documentclass[12pt,a4paper]{report}

% --- Фонт и језик (ћирилица) ---
\usepackage{fontspec}
\usepackage{polyglossia}
\setmainlanguage[script=cyrillic]{serbian}
\setotherlanguage{english}

\setmainfont{Times New Roman}
% Times New Roman из msttcorefonts има ћирилична слова, али не и OpenType ознаку
% писма, па се за ћирилицу користи исти фонт без те ознаке
\newfontfamily\cyrillicfont{Times New Roman}[Script=Default]
% Tahoma: наслови, насловна страна, апстракт, садржај, литература
\newfontfamily\tahoma{Tahoma}
% Моноспејс фонт са ћирилицом (за \texttt, нпр. називе команди)
\setmonofont{Liberation Mono}

% --- Математика: Times облик слова, као текст рада ---
\usepackage{amsmath}
\usepackage{unicode-math}
\setmathfont{TeX Gyre Termes Math}

% --- Маргине (шаблон: лева 3,1 cm; остале 2,5 cm) ---
% Висина текста и подножје одговарају шаблону: број стране је 2,7 cm од
% доње ивице, а последњи ред текста изнад подножја.
\usepackage{geometry}
\geometry{
	a4paper,
	left=3.1cm,
	right=2.5cm,
	top=2.5cm,
	textheight=672.2bp,
	footskip=21.2bp,
	nohead
}
\setlength{\topskip}{11.2bp}

% --- Величина слова 12 pt у Word-у је 12 bp (TeX-ов pt је мањи) ---
\makeatletter
\renewcommand\normalsize{\@setfontsize\normalsize{12bp}{14.5bp}}
\makeatother
\normalsize

% --- Проред 1,5 као у Word-у (1,5 × 13,8 pt = 20,7 pt); размак пасуса 12 pt ---
\usepackage{setspace}
\setstretch{1.42759}
\setlength{\parindent}{0bp}
\setlength{\parskip}{12bp}
% Шаблон не раставља речи на крају реда
\setlength{\emergencystretch}{3em}
\AtBeginDocument{\hyphenpenalty=10000 \exhyphenpenalty=10000}
% без растављања ред се развлачи (као у Word-у) уместо да излази на маргину
\tolerance=2000

% Висина великих слова Tahoma 14 pt (да наслов на врху стране стоји као у шаблону)
\newcommand{\naslovstrut}{\rule{0bp}{14bp}}
% (\strutbox: висина реда наслова као у Word-у, не 0,7 прореда)
\newcommand{\naslovfont}{\normalfont\tahoma\fontsize{14bp}{17.757bp}\selectfont
	\setbox\strutbox\hbox{\vrule height14bp depth2.89bp width0bp}}

% --- Набрајање: цртица (Tahoma), увлачење 0,63/1,26 cm; нумерисање „1.“ ---
\usepackage{enumitem}
\setlist{partopsep=0bp,parsep=0bp}
\setlist[itemize]{label={\tahoma -}, labelindent=17.85bp, labelwidth=17.85bp,
	labelsep=0bp, leftmargin=35.7bp, align=left,
	topsep=0bp, itemsep=1.2bp,
	before=\vspace{-5.2bp}, after=\vspace{0bp}}
\setlist[enumerate]{label=\arabic*., labelwidth=21.3bp, labelsep=18bp,
	leftmargin=39.3bp, align=right,
	topsep=0bp, itemsep=0bp}

% --- Наслови (Tahoma 14; поглавље не почиње нову страну) ---
% Размак изнад наслова се сабира са размаком изнад (као у Word-у, где се
% између два узастопна наслова узима већи од два размака)
\newcommand{\naslovpre}{\vskip-1.85bp\relax}
\usepackage{titlesec}
\titleclass{\chapter}{straight}
% 1. ниво: велика слова, затамњено
\titleformat{\chapter}[hang]
{\naslovpre\naslovfont\bfseries}
{\makebox[21.6bp][l]{\naslovstrut\thechapter}}{0bp}{\MakeUppercase}
% 2. ниво: велика слова
\titleformat{\section}[hang]
{\naslovpre\naslovfont}
{\makebox[28.8bp][l]{\naslovstrut\thesection}}{0bp}{\MakeUppercase}
% 3. ниво: мала слова
\titleformat{\subsection}[hang]
{\naslovpre\naslovfont}
{\makebox[36bp][l]{\naslovstrut\thesubsection}}{0bp}{}
\titlespacing*{\chapter}{0bp}{0bp}{1.8bp}
\titlespacing*{\section}{0bp}{0bp}{1.8bp}
\titlespacing*{\subsection}{0bp}{0bp}{1.8bp}

% --- Ненумерисани наслови (Литература): као у шаблону, увучени као нумерисани ---
\titleformat{name=\chapter,numberless}[hang]
{\naslovpre\naslovfont\bfseries}
{}{0bp}{\hspace*{21.6bp}\naslovstrut\MakeUppercase}

% --- Слике и табеле ---
\usepackage{graphicx}
\usepackage{caption}
\usepackage{float}
\usepackage{array}
\usepackage{chngcntr}
% Шаблон нумерише табеле, слике и формуле непрекидно (Табела 1, Слика 1, (1))
\counterwithout{table}{chapter}
\counterwithout{figure}{chapter}
\counterwithout{equation}{chapter}
% Натписи истим фонтом као текст; табеле: „Табела 1: …“ изнад табеле, лево;
% слике: „Слика 1. …“ испод слике, центрирано, ознака затамњена
\captionsetup{font={normalsize,stretch=1.42759},labelfont=normalfont,skip=0bp}
\captionsetup[table]{labelsep=colon,justification=raggedright,singlelinecheck=false,skip=10.6bp}
\captionsetup[figure]{labelsep=period,labelfont=bf,justification=centering,skip=2.7bp}
\setlength{\intextsep}{0bp}
% Табеле са пуним ивицама (0,5 pt), унутрашња маргина ћелије 0,19 cm;
% текст прве колоне поравнат са левом маргином
\setlength{\arrayrulewidth}{0.5bp}
\setlength{\tabcolsep}{5.4bp}
% Ћелије: 11 pt, проред 1,15 (у шаблону стил Table Grid даје 11 pt у табели)
\newcommand{\tabelalevo}{\setstretch{1}\fontsize{11bp}{14.5bp}\selectfont
	\hspace*{\dimexpr-\tabcolsep-\arrayrulewidth+0.4bp\relax}}
% Размак око табела и слика као око пасуса у шаблону
\BeforeBeginEnvironment{table}{\vspace{8.7bp}}
\AfterEndEnvironment{table}{\vspace{-10.1bp}}
\BeforeBeginEnvironment{figure}{\vspace{17.9bp}}
\AfterEndEnvironment{figure}{\vspace{-6.2bp}}

% --- Садржај (као у шаблону: без тачкица и затамњења, проред 1,5) ---
\usepackage{tocloft}
\tocloftpagestyle{empty}
\renewcommand{\cfttoctitlefont}{\naslovfont\bfseries\naslovstrut}
\setlength{\cftbeforetoctitleskip}{-3.3bp}
\setlength{\cftaftertoctitleskip}{34.5bp}
\renewcommand{\cftchapfont}{\normalfont}
\renewcommand{\cftchappagefont}{\normalfont}
\renewcommand{\cftchapleader}{\hfill}
\renewcommand{\cftsecleader}{\hfill}
\renewcommand{\cftsubsecleader}{\hfill}
\setlength{\cftbeforechapskip}{12.15bp}
\setlength{\cftbeforesecskip}{12.15bp}
\setlength{\cftbeforesubsecskip}{12.15bp}
\setlength{\cftchapindent}{0bp}
\setlength{\cftchapnumwidth}{22bp}
\setlength{\cftsecindent}{11bp}
\setlength{\cftsecnumwidth}{33bp}
\setlength{\cftsubsecindent}{22bp}
\setlength{\cftsubsecnumwidth}{44bp}
\addto\captionsserbian{%
	\renewcommand{\listfigurename}{ЛИСТА СЛИКА}%
	\renewcommand{\listtablename}{ЛИСТА ТАБЕЛА}%
	\renewcommand{\contentsname}{САДРЖАЈ}%
}

% --- Формуле: размак као између пасуса ---
\makeatletter
\g@addto@macro\normalsize{%
	\setlength{\abovedisplayskip}{17.2bp}%
	\setlength{\belowdisplayskip}{17.2bp}%
	\setlength{\abovedisplayshortskip}{17.2bp}%
	\setlength{\belowdisplayshortskip}{17.2bp}%
}
\makeatother

% --- Литература: APA стил ---
\usepackage[
style=apa,
backend=biber,
language=serbian,
uniquelist=false, % шаблон: више од два аутора → „први аутор и др“
autolang=hyphen % енглеске референце се растављају по енглеским правилима
]{biblatex}
\addbibresource{references.bib}
% Изглед списка литературе као текст рада: обострано, проред 1,5,
% без висећег увлачења, 12 pt између ставки
\renewcommand*{\bibfont}{\normalsize}
\defbibenvironment{bibliography}
{\list{}{%
		\setlength{\leftmargin}{0bp}%
		\setlength{\itemindent}{0bp}%
		\setlength{\labelwidth}{0bp}%
		\setlength{\labelsep}{0bp}%
		\setlength{\topsep}{1.8bp}%
		\setlength{\partopsep}{0bp}%
		\setlength{\itemsep}{12bp}%
		\setlength{\parsep}{0bp}}%
	\renewcommand*{\makelabel}[1]{##1\hss}}
{\endlist}
{\item}
% Шаблон: „(Јовић и др, 2021)“
\DefineBibliographyStrings{serbianc}{andothers={и\addabbrvspace др}}
% Линкови се ломе на било ком знаку, без развлачења
\setcounter{biburlnumpenalty}{100}
\setcounter{biburlucpenalty}{100}
\setcounter{biburllcpenalty}{100}

\usepackage{tikz}

\usepackage[hidelinks]{hyperref}


% --- Број стране десно у подножју (као у шаблону) ---
\usepackage{fancyhdr}
\fancypagestyle{radstil}{%
	\fancyhf{}%
	\renewcommand{\headrulewidth}{0bp}%
	\fancyfoot[R]{\thepage}%
}
\fancypagestyle{plain}{%
	\fancyhf{}%
	\renewcommand{\headrulewidth}{0bp}%
	\fancyfoot[R]{\thepage}%
}

% =============================================================================
\begin{document}

	% Насловна страна, посвета, апстракт и садржај имају у шаблону маргине
	% 2,54 cm (доња 1,5 cm)
	\newgeometry{left=72bp,right=72bp,top=72bp,textheight=717.3bp,nohead,nofoot}

	% -----------------------------------------------------------------------------
	% НАСЛОВНА СТРАНА (према шаблону: Прилог 2а)
	% Положај сваког реда је измерен у шаблону (у pt од горње ивице стране).
	% -----------------------------------------------------------------------------
	\pagestyle{empty}
	\pagenumbering{gobble} % предњи део се не нумерише
	\begin{titlepage}
		\thispagestyle{empty}
		\begin{tikzpicture}[remember picture,overlay,
			every node/.style={inner sep=0bp,outer sep=0bp}]
			\tahoma\bfseries
			\node[anchor=base] at ([yshift=-98.0bp]current page.north)
			{\fontsize{14bp}{14bp}\selectfont УНИВЕРЗИТЕТ У БЕОГРАДУ};
			\node[anchor=base] at ([yshift=-126.9bp]current page.north)
			{\fontsize{14bp}{14bp}\selectfont ФАКУЛТЕТ ОРГАНИЗАЦИОНИХ НАУКА};

			\node[anchor=base] at ([yshift=-302.8bp]current page.north)
			{\fontsize{14bp}{14bp}\selectfont ЗАВРШНИ РАД};
			% наслов: Tahoma 20, проред 1,5 (36,2 pt); преломи као у Word-у
			% (ред се ломи када наредна реч не стане у ширину текста)
			\foreach \red/\tekst in {0/{Утицај скупа алата на успешност и трошак},
				1/{агената заснованих на великим језичким},
				2/{моделима при решавању софтверских},
				3/{проблема}}
			\node[anchor=base] at ([yshift={-346.2bp-\red*36.2bp}]current page.north)
			{\fontsize{20bp}{20bp}\selectfont\tekst};

			\node[anchor=base west] at ([xshift=77.5bp,yshift=-541.3bp]current page.north west)
			{\fontsize{14bp}{14bp}\selectfont Ментор};
			\node[anchor=base west] at ([xshift=318.6bp,yshift=-541.3bp]current page.north west)
			{\fontsize{14bp}{14bp}\selectfont Студент};
			\node[anchor=base west,text width=230.3bp,align=left,
			font=\fontsize{14bp}{11.838bp}\selectfont]
			at ([xshift=77.5bp,yshift=-575.3bp]current page.north west)
			{Др Андрија Петровић, Доцент};
			\node[anchor=base west,text width=199.2bp,align=left,
			font=\fontsize{14bp}{17.757bp}\selectfont]
			at ([xshift=318.6bp,yshift=-575.3bp]current page.north west)
			{Алекса Љујић, 2021/0276};

			\node[anchor=base] at ([yshift=-745.3bp]current page.north)
			{\fontsize{14bp}{14bp}\selectfont Београд, 2026. године};
		\end{tikzpicture}
		\null
	\end{titlepage}

	% -----------------------------------------------------------------------------
	% ПОСВЕТА (по избору — обрисати ако се не користи)
	% -----------------------------------------------------------------------------
	\newpage
	\thispagestyle{empty}
	\begin{flushright}
		\itshape Посвета (по избору)
	\end{flushright}

	% -----------------------------------------------------------------------------
	% АПСТРАКТ (око 200 речи)
	% -----------------------------------------------------------------------------
	\newpage
	\thispagestyle{empty}
	{\naslovfont\bfseries\rule{0bp}{26bp}АПСТРАКТ\par}
	\vspace{1.8bp}

	Развој великих језичких модела довео је до појаве агентских система који, поред генерисања текста и кода, могу самостално да користе алате и делују у софтверском окружењу. У оваквим системима скуп доступних алата представља важан део интерфејса између језичког модела и окружења, али није у потпуности јасно у којој мери појединачни алати доприносе успешности решавања задатака и да ли њихова употреба оправдава додатни трошак. Циљ овог рада је испитивање утицаја скупа алата на успешност, трошак и способност локализације грешке код агента за решавање проблема у софтверском инжењерству, при фиксираном језичком моделу. У ту сврху развијен је сопствени агентски систем заснован на петљи позивања алата, са извршавањем команди у изолованим Docker контејнерима. Експеримент је спроведен на свих 300 задатака из скупа SWE-bench Lite, у факторском дизајну $2 \times 2$, којим се независно испитује утицај алата за детерминистичку замену текста и групе алата за претрагу и читање датотека, док је алат за извршавање команди доступан у свим условима. Као примарне мере посматрају се стопа решених задатака, трошак по решеном задатку и стопа локализације грешке, док се додатно анализирају потрошња токена, број корака, време извршавања и коришћење појединачних алата. Поређење конфигурација заснива се на упареним статистичким поступцима, интервалима поверења и анализи неинфериорности. Посебна пажња посвећена је разликовању ефекта доступности алата од ефекта њиховог стварног коришћења. Комбинација љуске и алата за замену текста решила је 119 од 300 задатака (39,7\%), а пун скуп алата 84 (28,0\%), чиме је потврђена неинфериорност мањег скупа, уз интервал поверења разлике који у целини лежи изнад нуле. Алати за претрагу повећали су трошак по покретању, али нису побољшали локализацију грешке, већ су је смањили за 4,7 процентних поена. Хипотеза да алат за замену текста смањује трошак по решеном задатку није потврђена, јер трошак агента претежно одређују улазни токени, а не величина измене. Закључује се да код малог модела већи скуп алата не даје бољег агента, а да најбољи однос успешности и трошка међу испитаним условима пружа минималан скуп који садржи љуску и наменски алат за измену датотека.
	
	\textbf{Кључне речи:} велики језички модели, агенти, агентски системи, позивање алата, софтверско инжењерство, SWE-bench Lite, аблациона анализа, локализација грешке, трошак, однос цене и перформанси.
	
	% -----------------------------------------------------------------------------
	% САДРЖАЈ
	% -----------------------------------------------------------------------------
	\newpage
	\thispagestyle{empty}
	\tableofcontents
	
	% -----------------------------------------------------------------------------
	% НАПОМЕНА: шаблон (Прилог 2а) не предвиђа листу акронима, листу слика и
	% листу табела. Ако их ментор тражи, откоментарисати наредни блок.
	% -----------------------------------------------------------------------------
	% \newpage
	% \phantomsection
	% \addcontentsline{toc}{chapter}{Листа акронима}
	% \chapter*{Листа акронима}
	% \begin{description}
	%	\item[ACI] интерфејс између агента и рачунара (енг. \textit{Agent-Computer Interface})
	%	\item[API] интерфејс за програмирање апликација (енг. \textit{Application Programming Interface})
	%	\item[CSV] вредности раздвојене зарезом (енг. \textit{Comma-Separated Values})
	%	\item[JSON] запис објеката у JavaScript-у (енг. \textit{JavaScript Object Notation})
	%	\item[LLM] велики језички модел (енг. \textit{Large Language Model})
	%	\item[SWE] софтверско инжењерство (енг. \textit{Software Engineering})
	% \end{description}
	% \newpage
	% \listoffigures
	% \newpage
	% \listoftables
	
	% =============================================================================
	% ПОГЛАВЉА РАДА
	% =============================================================================
	\restoregeometry
	\pagestyle{radstil}
	\pagenumbering{arabic}
	\setcounter{page}{1}
	
	% =============================== 1. УВОД =====================================
	% у шаблону увод почиње после прелома одељка, па задржава размак од 12 pt изнад
	\vspace*{-20.7bp}
	\chapter{Увод}
	
	Велики језички модели су у софтверском инжењерству најпре коришћени за
	допуњавање кода у развојном окружењу и као асистенти у облику разговора. У таквом начину рада контекст држи човек. Он бира који део кода
	ће моделу показати, преноси предлог решења у пројекат, покреће програм и
	проверава да ли измена заиста решава проблем. Модел даје предлог, али не
	делује у окружењу.

	Последњих година тежиште се помера ка вештачкој интелигенцији као агенту.
	Агент је систем у коме језички модел у петљи сам бира наредни корак: чита
	датотеке репозиторијума, покреће команде, мења код и проверава резултат,
	све док задатак не буде решен. Истраживачки системи као што су SWE-agent
	\parencite{yang2024sweagent} и OpenHands \parencite{wang2025openhands}, као
	и комерцијални алати засновани на истом принципу, показују да агенти могу
	да решавају стварне проблеме из репозиторијума отвореног кода. Скуп података
	SWE-bench \parencite{jimenez2024swebench} постао је стандардна мера
	напретка у овој области.

	Успешност агента, међутим, не зависи само од модела. Систем који окружује
	модел (енг. \textit{harness}) одређује које алате модел има на располагању,
	како су описани и какве повратне информације добија. \textcite{yang2024sweagent}
	показују да дизајн интерфејса између агента и рачунара значајно утиче на
	резултат, иако је модел исти. Сваки алат ипак има и цену: његов опис се шаље
	моделу у сваком позиву, а резултати позива продужавају контекст, који се у
	наредним корацима поново обрађује и плаћа. Истовремено,
	\textcite{xia2025agentless} показују да се конкурентни резултати могу
	постићи и знатно једноставнијим системом, што доводи у питање да ли је
	сложенији скуп алата заиста потребан. Недостаје систематско мерење
	изолованог доприноса појединачних алата, и то не само на успешност, него и
	на трошак.

	Питање трошка има практичну тежину. Коришћење модела путем API-ја плаћа се
	по токену, а агент у једном задатку прави десетине позива модела. Мали и
	јефтини модели су посебно занимљиви за примену, али је управо код њих
	отворено колико добро користе алате које им систем пружа. Овај рад се бави
	тим питањем на примеру малог модела и стварних задатака из SWE-bench скупа.

	\section{Циљ рада}

	Општи циљ рада је да се измери како скуп алата доступних агенту утиче на
	успешност решавања задатака из софтверског инжењерства и на трошак тог
	решавања, при фиксираном језичком моделу.

	Конкретни циљеви рада су:
	\begin{itemize}
		\item имплементирати сопствени агентски систем са скупом алата који се
		може конфигурисати и изолованим окружењем за извршавање;
		\item спровести аблацију скупа алата у факторском дизајну $2 \times 2$:
		присуство алата за замену текста (\texttt{str\_replace}) и присуство
		групе алата за претрагу и читање датотека, уз љуску (\texttt{bash}) као
		основни алат у свим условима;
		\item измерити стопу решених задатака званичном SWE-bench процедуром
		евалуације, трошак по решеном задатку, као и способност агента да
		пронађе датотеку у којој се налази грешка;
		\item утврдити која конфигурација алата пружа најбољи однос цене и
		перформанси, а која доноси предност при лоцирању грешке;
		\item утврдити у којој мери агент заиста користи алате који су му
		доступни, како би се ефекат доступности алата разликовао од ефекта
		његовог коришћења.
	\end{itemize}

	\section{Методологија}

	За потребе рада развијен је сопствени агентски систем у програмском језику
	Python, чији је изворни код јавно доступан у репозиторијуму
	\url{https://github.com/aleksaljujic/agent-harness}. Агент ради у петљи позивања алата, а свака команда се извршава у
	изолованом Docker контејнеру без приступа мрежи, тако да агент не може да
	дође до званичне исправке проблема. Као модел се користи \texttt{gpt-5.4-nano},
	са температуром 0 и искљученим режимом проширеног закључивања. Модел је
	намерно фиксиран, јер је предмет рада скуп алата, а не поређење модела.

	Експерименту је претходила пилот евалуација на мањем, синтетички
	направљеном скупу задатака опште намене (креирање и измена датотека,
	исправљање грешака, рад са тестовима). У пилоту је комбинација
	\texttt{bash} и \texttt{str\_replace} остварила најнижи трошак и најкраће
	време извршавања уз високу успешност. Пошто је пилот рађен на малом броју
	задатака, са по једним покретањем и на вештачки направљеним задацима, тај
	налаз није довољан за закључак. Он је, међутим, мотивисао систематски
	експеримент на стварним задацима, описан у наставку.

	Главни експеримент спроводи се на целом скупу SWE-bench Lite (300 задатака),
	у четири конфигурације алата, што даје 1200 покретања. Све конфигурације
	користе идентичан системски упит, који не помиње ниједан конкретан алат,
	тако да се између услова разликује само скуп доступних алата. Решења се
	оцењују званичним SWE-bench тестовима. Примарне мере су стопа решених
	задатака, трошак по решеном задатку и стопа локализације, тј. удео
	покретања у којима исправка агента мења бар једну датотеку коју мења и
	референтна исправка. Секундарне мере су трошак по покретању, број токена,
	број корака и време извршавања. Пошто сваки задатак пролази кроз све
	конфигурације, поређења су упарена. Мере, правила искључивања и
	конфигурација записани су пре покретања експеримента, а хипотезе и план
	статистичке анализе пре оцењивања решења и пре рачунања стопе
	локализације. Свако покретање бележи верзију кода којом је добијено. Детаљи су описани у
	поглављу~4.

	\section{Хипотезе}

	Полазна теза рада је да комбинација љуске и алата за замену текста
	(\texttt{bash} и \texttt{str\_replace}) пружа најбољи однос цене и
	перформанси, док пун скуп алата, који укључује и алате за претрагу и читање
	датотека, доноси предност пре свега при лоцирању грешке. Да би се ова теза
	могла проверити, рашчлањена је на три хипотезе:

	\begin{itemize}
		\item \textbf{Х1:} Доступност алата за детерминистичку замену текста
		(\texttt{str\_replace}) смањује трошак по решеном задатку;
		\item \textbf{Х2:} Доступност алата за претрагу и читање датотека
		повећава трошак по покретању, али повећава стопу локализације, тј. удео
		покретања у којима агент измени датотеку у којој се налази грешка;
		\item \textbf{Х3:} Стопа решених задатака конфигурације са алатима
		\texttt{bash} и \texttt{str\_replace} није нижа од стопе конфигурације са
		пуним скупом алата за више од пет процентних поена.
	\end{itemize}

	Хипотезе Х1 и Х2 проверавају се на главним ефектима факторског дизајна:
	ефекат алата \texttt{str\_replace} добија се поређењем свих конфигурација са
	тим алатом и без њега, при оба нивоа другог фактора, а ефекат алата за
	претрагу и читање на исти начин. Хипотеза Х3 је хипотеза неинфериорности и
	служи као контрола да уштеда из Х1 није постигнута на рачун лошијих
	резултата. Ако се потврде све три хипотезе, из њих следи полазна теза:
	конфигурација са алатима \texttt{bash} и \texttt{str\_replace} је
	најисплативији избор, а алати за претрагу и читање имају вредност при
	лоцирању грешке, али уз већи трошак. Разлике у ефекту између појединих
	врста задатака анализирају се експлоративно, без унапред формулисане
	хипотезе.

	\section{Структура рада}

	Рад је организован у шест поглавља. У другом поглављу дат је преглед
	литературе о агентима који користе алате, бенчмарковима за евалуацију
	генерисања кода и агентским системима за решавање проблема у
	репозиторијумима. Треће поглавље уводи теоријске основе: велике језичке
	моделе, појам агента и позивања алата, специфичности агената који раде над
	кодом и евалуациони оквир. Четврто поглавље описује методологију:
	архитектуру развијеног агента, дизајн алата и конфигурације које се пореде,
	окружење за извршавање и поступак евалуације. У петом поглављу приказани су
	поставка и резултати експеримента, њихова дискусија у односу на хипотезе и
	ограничења рада. Шесто поглавље сумира главне налазе и наводи правце даљег
	истраживања.
	
	% ======================== 2. РЕЛЕВАНТНА ЛИТЕРАТУРА ===========================
	\chapter{Релевантна литература}
	
	% НАЦРТ: тематски преглед по групама литературе. Проширити након читања
	% радова (препоручени редослед: Agentless, SWE-agent, SWE-Bench+, Card et al.).

	Преглед литературе организован је тематски: најпре су приказани темељни
	радови о агентима који користе алате, затим бенчмаркови за евалуацију
	генерисања кода и њихова ограничења, агентски системи за решавање проблема
	у стварним репозиторијумима, радови о коришћењу алата и дугачког контекста,
	и на крају методолошки радови о статистички поузданој евалуацији.

	\section{Темељи агената који користе алате}

	Парадигма ReAct \parencite{yao2023react} уводи наизменично смењивање корака
	закључивања и деловања, чиме језички модел може да планира, позива спољашње
	алате и прилагођава план на основу добијених опсервација. Ова петља
	представља основу већине савремених агентских система, укључујући и систем
	развијен у овом раду. Паралелно, \textcite{schick2023toolformer} показују да
	модел може самостално да научи када и како да позове алат (енг.
	\textit{tool learning}), чиме се позивање алата успоставља као способност
	модела, а не само као спољашња надоградња.

	\section{Бенчмаркови за генерисање кода и квалитет скупова података}

	Рана евалуација језичких модела над кодом ослањала се на кратке, самосталне
	задатке, као што су HumanEval \parencite{chen2021humaneval} и MBPP
	\parencite{austin2021mbpp}. Каснији бенчмаркови проширују обухват ка
	сложенијим позивима библиотека \parencite{zhuo2024bigcodebench} и ка
	евалуацији отпорној на контаминацију тренинг података
	\parencite{jain2024livecodebench}. \textcite{liu2023evalplus} показују да
	слаби тестови систематски прецењују перформансе модела, што је директно
	релевантно за метрику коришћену у овом раду.

	SWE-bench \parencite{jimenez2024swebench} помера евалуацију ка стварним
	проблемима (енг. \textit{issues}) из GitHub репозиторијума, где модел мора
	да разуме постојећу кодну базу и произведе исправку (енг. \textit{patch}), тј.
	скуп измена изворног кода који пролази задате тестове.
	Овај скуп података и његова евалуациона процедура представљају основу
	експерименталног дела рада. Проширења обухватају визуелне домене
	\parencite{yang2024swebenchmm} и више програмских језика
	\parencite{zan2025multiswebench}. Истовремено,
	\textcite{aleithan2024swebenchplus} документују озбиљне проблеме квалитета
	у SWE-bench-у: у 32,67\% успешних исправки решење је било садржано у самом
	опису проблема или коментарима (енг. \textit{solution leakage}), а 31,08\%
	прихваћених исправки сумњиво је због тестова који нису довољни да провере
	исправност. Након уклањања проблематичних задатака, стопа решених задатака
	система SWE-agent са моделом GPT-4 пала је са 12,47\% на 3,97\%. Ови налази
	се разматрају у поглављу о ограничењима.

	\section{Агентски системи за решавање проблема у софтверском инжењерству}

	\textcite{yang2024sweagent} уводе појам интерфејса између агента и рачунара
	(енг. \textit{Agent-Computer Interface}, ACI) и показују да дизајн скупа
	алата значајно утиче на успешност агента. Аблацијом на SWE-bench Lite скупу
	пореде различите начине претраге, прегледа и измене датотека. Посебно је
	релевантно поређење са агентом који има само љуску: исти модел решава
	18,00\% задатака са прилагођеним интерфејсом, а 11,00\% само са љуском. Ова
	аблација представља најближи претходни рад за хипотезе о утицају скупа
	алата. OpenHands \parencite{wang2025openhands} је отворена платформа
	опште намене за софтверске агенте; у овом раду је ипак развијен сопствени
	систем, како би се обезбедила потпуна контрола над скупом алата који се
	пореди. HyperAgent \parencite{phan2024hyperagent} је систем са четири
	специјализована агента (планер, навигатор, уређивач кода и извршилац). Аутори
	спроводе аблацију улога агената и анализу утицаја појединих алата, а уз
	стопу решених задатака извештавају и просечан трошак по задатку, што га
	чини методолошки најближим овом раду. Међу осталим приступима
	издвајају се AutoCodeRover \parencite{zhang2024autocoderover}, који
	комбинује агента са техникама анализе програма, као и оквири за више
	агената попут AutoGen-а \parencite{wu2023autogen}.

	Посебно је важан рад \textcite{xia2025agentless}, који доводи у питање
	потребу за сложеним агентима. Њихов приступ Agentless, заснован на
	трофазном процесу без агентске петље, постиже 32,00\% (96 тачних исправки)
	на SWE-bench Lite скупу уз просечан трошак од 0,70 долара, надмашујући
	тадашње агенте отвореног кода. Аутори такође конструишу SWE-bench Lite-S
	уклањањем проблематичних инстанци \parencite{xia2025agentless}. Овај рад полази од сродног питања, али
	из другог угла: уместо да уклони агентску петљу, задржава је и мери да ли
	специјализовани алати доносе стварну корист у односу на рад само са љуском.

	Други правац истраживања фокусира се на тренирање модела за решавање
	проблема у репозиторијумима: SWE-Fixer \parencite{xie2025swefixer},
	SWE-Gym \parencite{pan2024swegym}, SWE-RL \parencite{wei2025swerl} и
	SWE-smith \parencite{yang2025swesmith}. Ти приступи мењају сам модел, кроз
	додатно обучавање или учење поткрепљивањем, па захтевају податке за обучавање
	и знатне рачунарске ресурсе, а резултат важи за модел који је на тај начин
	добијен. Овај рад полази са друге стране: модел остаје непромењен, а мења се
	само оно што га окружује, односно скуп алата. Два приступа су допунска, јер
	се налази о дизајну алата могу применити на било који модел, укључујући и
	оне посебно обучене за решавање проблема у репозиторијумима.

	\section{Коришћење алата, контекст и ефикасност}

	ToolLLM \parencite{qin2023toolllm} и Gorilla \parencite{patil2023gorilla}
	проучавају способност модела да ефикасно користе велики број API-ја, док
	$\tau$-bench \parencite{yao2024taubench} евалуира агенте у интеракцији са
	алатима и корисником. \textcite{liu2024lostinthemiddle} показују да модели
	слабије користе информације смештене у средину дугачког контекста, што
	објашњава зашто више контекста не значи нужно боље резултате и подржава
	претпоставку да мањи скуп алата може бити ефикаснији. У истом духу,
	\textcite{xiao2025trajectory} уклањају сувишне информације из трајекторије
	агента и тиме смањују број улазних токена за 39,9--59,7\%, а укупан трошак
	за 21,1--35,9\%, уз очуване перформансе.

	\section{Методологија евалуације и статистичка поузданост}

	\textcite{bowman2021fix} разматрају системске проблеме бенчмаркова у обради
	природног језика. \textcite{card2020littlepower} показују да је велики број
	експеримената у овој области статистички подмоћан и дају смернице за
	процену потребне величине узорка. \textcite{dror2018hitchhiker} пружају
	практичан водич за избор статистичког теста, укључујући упарене тестове, а
	\textcite{reimers2017distributions} образлажу зашто треба извештавати
	дистрибуцију резултата уместо појединачне вредности. Ови радови усмеравају
	статистичку обраду резултата у поглављу~4.

	\section{Празнина у литератури}

	Претходни радови показују да дизајн скупа алата утиче на успешност агента,
	али то питање углавном обрађују на један од два начина. Први је поређење
	целих система, као у радовима SWE-agent \parencite{yang2024sweagent},
	OpenHands \parencite{wang2025openhands} и Agentless
	\parencite{xia2025agentless}, где се истовремено разликују и алати, и
	упити, и организација рада агента, па се допринос појединачних алата не
	може издвојити. Други је аблација у којој се компоненте уклањају једна по
	једна, као у радовима \textcite{yang2024sweagent} и
	\textcite{phan2024hyperagent}, што показује да ли компонента помаже, али не
	и како се ефекти различитих компоненти међусобно односе. Факторски нацрт, у
	коме се две компоненте укрштају и посматра се и њихова интеракција,
	примењиван је и на друге делове агентског система, на пример на планирање и
	памћење ранијих покушаја \parencite{zhang2026pmcoder}, али не и на скуп
	алата. Поред тога,
	евалуације се најчешће спроводе са најјачим доступним моделима, док је
	мање познато колико успешно мали и јефтини модели решавају стварне задатке
	и колико им при томе помажу алати.

	Овај рад доприноси у два правца. Главни допринос је примена факторског
	нацрта на скуп алата агента, односно систематско мерење ефекта две
	функционално различите групе: алата за измену кода
	(\texttt{str\_replace}) и алата за претрагу и читање датотека. Групе се
	посматрају као два независна фактора у факторском дизајну $2 \times 2$,
	уз љуску као заједничку основу, при чему системски упит и сви остали
	услови остају исти. Тиме се ефекат сваке групе мери одвојено, при оба нивоа
	друге групе, а могуће је уочити и њихову међусобну интеракцију. Ефекат се
	не мери само на стопи решених задатака, већ и на трошку по решеном задатку
	и на способности агента да пронађе датотеку у којој се налази грешка.

	Споредни допринос рада је увид у то колико успешно мали модел решава
	стварне проблеме из софтверског инжењерства. Експеримент се у целини
	спроводи са моделом \texttt{gpt-5.4-nano}, при температури 0 и без режима
	проширеног закључивања, на свих 300 задатака из скупа SWE-bench Lite.
	Резултати тако дају и основну меру могућности малог модела у агентском
	окружењу, као и одговор на питање колико му у томе помажу поједине групе
	алата.

	% ========================== 3. ТЕОРИЈСКЕ ОСНОВЕ ==============================
	\chapter{Теоријске основе}
	
	Ово поглавље уводи појмове на којима почива остатак рада. Најпре су описани
	велики језички модели као основа система, затим појам агента и механизам
	позивања алата, који моделу омогућава да делује у окружењу. Након тога се
	разматрају специфичности агената који раде над изворним кодом, а поглавље
	се завршава описом евалуационог оквира: скупа података, метрика успешности
	и трошка, као и статистичких поступака за њихово поређење.

	\section{Велики језички модели}

	\subsection{Језичко моделовање}

	Језички модел дефинише расподелу вероватноћа над низовима токена, где је
	токен основна јединица текста (реч, део речи или симбол). Савремени модели
	су ауторегресивни: вероватноћа низа $x = (x_1, \dots, x_T)$ разлаже се на
	производ условних вероватноћа сваког токена у односу на претходне:
	\begin{equation}
		p_\theta(x) = \prod_{t=1}^{T} p_\theta(x_t \mid x_{<t}),
		\label{eq:autoregresija}
	\end{equation}
	где $\theta$ означава параметре модела. Модел се тренира минимизацијом
	негативне лог-веродостојности над великим корпусом текста:
	\begin{equation}
		\mathcal{L}(\theta) = -\sum_{t=1}^{T} \log p_\theta(x_t \mid x_{<t}).
	\end{equation}
	Дакле, основни задатак модела јесте предикција следећег токена. Сложеније
	способности, попут превођења, закључивања или писања кода, настају као
	последица тог једноставног циља примењеног на довољно велики модел и
	довољно велику количину података.

	\subsection{Трансформер архитектура}

	Готово сви савремени велики језички модели (енг. \textit{Large Language
	Models}, LLM) засновани су на трансформер архитектури
	\parencite{vaswani2017attention}. Њен централни механизам је пажња (енг.
	\textit{attention}), која сваком токену омогућава да директно приступи
	репрезентацијама свих претходних токена у контексту:
	\begin{equation}
		\mathrm{Attention}(Q, K, V) = \mathrm{softmax}\!\left(\frac{QK^\top}{\sqrt{d_k}}\right) V,
	\end{equation}
	где су $Q$, $K$ и $V$ матрице упита, кључева и вредности добијене линеарним
	пројекцијама улаза, а $d_k$ димензија кључева. Пошто се пажња рачуна између
	свих парова токена, рачунска и меморијска сложеност расту квадратно са
	дужином контекста. Због тога модели имају ограничен контекстни прозор, а
	обрада дугачких улаза је скупља и спорија.

	\subsection{Генерисање и понашање модела у пракси}

	При генерисању, модел у сваком кораку израчунава логите $z$ за све токене
	речника, који се претварају у вероватноће уз параметар температуре $\tau$:
	\begin{equation}
		p(x_t = i \mid x_{<t}) = \frac{\exp(z_i / \tau)}{\sum_j \exp(z_j / \tau)}.
	\end{equation}
	Ниже вредности $\tau$ чине излаз детерминистичнијим, а више разноврснијим.
	Чак и при $\tau = 0$, излаз комерцијалних модела у пракси није потпуно
	детерминистички, што има директне последице по поновљивост експеримената.

	\textcite{brown2020gpt3} показују да довољно велики модели могу да решавају
	нове задатке само на основу инструкција и примера у улазу (енг.
	\textit{in-context learning}), без промене параметара. Додатно подешавање
	модела да прати инструкције, укључујући учење на основу људских повратних
	информација \parencite{ouyang2022instructgpt}, чини моделе употребљивим као
	асистенте опште намене. Навођење модела да пре коначног одговора експлицитно
	изложи међукораке закључивања (енг. \textit{chain-of-thought})
	значајно побољшава резултате на задацима који захтевају више корака
	\parencite{wei2022cot}.

	\subsection{Ограничења}

	Велики језички модели имају неколико ограничења која су релевантна за овај
	рад. Прво, знање модела је ограничено на податке из тренинга и не обухвата
	конкретну кодну базу над којом се ради, па модел мора да прикупи информације
	из окружења. Друго, модели могу да генеришу садржај који звучи уверљиво, али
	је нетачан (енг. \textit{hallucination}) \parencite{ji2023hallucination},
	нпр. да наведу непостојећу функцију или датотеку. Треће, иако контекстни
	прозори савремених модела обухватају стотине хиљада токена, модели не
	користе подједнако добро све делове контекста:
	\textcite{liu2024lostinthemiddle} показују да се информације у средини
	дугачког улаза користе знатно слабије од оних на почетку и крају. Четврто,
	трошак коришћења модела путем API-ја обрачунава се по броју улазних и
	излазних токена, па свака додатна информација у контексту има директну цену.

	\section{Агенти}

	\subsection{Појам агента}

	У класичној литератури вештачке интелигенције, агент је систем који
	опажа своје окружење путем сензора и делује на њега путем актуатора, са
	циљем да оствари задати циљ \parencite{russell2020aima}. У контексту
	великих језичких модела, агент је систем у коме језички модел понавља
	циклус: на основу тренутног стања одлучује коју акцију да предузме, акција
	се извршава у окружењу, а резултат се враћа моделу као нова опсервација.
	За разлику од једноставног упита, где модел даје један одговор, агент
	самостално одређује број и редослед корака потребних за решавање задатка.

	\subsection{Позивање алата}

	Језички модел сам по себи производи само текст. Да би деловао у окружењу,
	користи алате (енг. \textit{tools}): функције које извршава спољашњи
	систем, као што су читање датотеке, претрага или покретање команде.
	\textcite{schick2023toolformer} показују да модел може да научи када и са
	којим аргументима да позове алат, а радови попут ToolLLM
	\parencite{qin2023toolllm} и Gorilla \parencite{patil2023gorilla}
	проширују ову способност на велики број стварних API-ја. Пошто је реч о
	наученој способности, на понашање модела утиче и то колико су назив и облик
	понуђеног алата слични онима на којима је модел обучаван.

	У савременим системима, сваки алат се описује називом, описом на природном
	језику и шемом параметара (најчешће у JSON Schema формату). Опис свих
	доступних алата укључује се у контекст модела, а модел, уместо слободног
	текста, може да генерише структурисани позив алата са аргументима. Систем
	који покреће агента (енг. \textit{harness}) извршава позив и резултат враћа
	моделу. Важна последица је да \emph{сваки алат заузима део контекста}:
	описи алата се шаљу у сваком позиву модела, а већи број алата повећава и
	број могућих избора које модел мора да разликује.

	\subsection{Петља закључивања и деловања}

	Парадигма ReAct \parencite{yao2023react} комбинује закључивање и деловање:
	модел наизменично генерише мисао (образложење следећег корака) и акцију
	(позив алата), након чега добија опсервацију. Формално, нека је $q$ опис
	задатка. У кораку $t$ контекст модела је
	\begin{equation}
		c_t = (q, r_1, a_1, o_1, \dots, r_{t-1}, a_{t-1}, o_{t-1}),
	\end{equation}
	где су $r_i$ мисли, $a_i$ акције, а $o_i$ опсервације. Модел генерише
	наредну мисао и акцију $(r_t, a_t) \sim p_\theta(\cdot \mid c_t)$, окружење
	враћа опсервацију $o_t = \mathrm{Env}(a_t)$, и петља се наставља док модел
	не сигнализира крај или док се не достигне ограничење броја корака или
	трошка. Низ $(r_1, a_1, o_1, \dots, r_n, a_n, o_n)$ назива се трајекторија
	агента.

	Из ове формулације следи важно својство трошка. Пошто се у сваком кораку
	моделу шаље цео досадашњи контекст, укупан број улазних токена у
	трајекторији дужине $n$ износи
	\begin{equation}
		N_{\text{in}} = \sum_{t=1}^{n} |c_t|,
		\label{eq:ulazni_tokeni}
	\end{equation}
	што расте приближно квадратно са бројем корака ако свака опсервација
	додаје сличан број токена. Опширне опсервације (нпр. цео садржај великих
	датотека или дугачак излаз команди) и дугачки описи алата зато
	вишеструко повећавају трошак. Кеширање префикса упита (енг.
	\textit{prompt caching}) код појединих провајдера смањује цену поновљених
	делова контекста, али не и њихов утицај на дужину контекста.
	\textcite{xiao2025trajectory} показују да се скраћивањем трајекторија може
	значајно смањити трошак агентских система.

	\subsection{Архитектуре агентских система}

	Агентски системи се разликују по степену аутономије. На једном крају су
	фиксни процеси у којима програмер унапред одређује кораке, а модел решава
	појединачне подзадатке; пример је Agentless \parencite{xia2025agentless},
	који задатак дели на локализацију, поправку и валидацију. На другом крају
	су аутономни агенти, у којима модел сам бира наредну акцију, као у
	SWE-agent \parencite{yang2024sweagent} и OpenHands
	\parencite{wang2025openhands}. Постоје и системи са више агената који
	имају различите улоге и међусобно комуницирају, као AutoGen
	\parencite{wu2023autogen} и HyperAgent \parencite{phan2024hyperagent}.
	Већа аутономија пружа флексибилност, али повећава број позива модела,
	трошак и варијабилност резултата.

	\section{Агенти за извршавање кода}

	\subsection{Окружење и акције}

	Агенти за софтверско инжењерство раде над стварним репозиторијумом кода.
	Њихово окружење чине систем датотека, љуска (енг. \textit{shell}) и
	програмски алати пројекта, као што су интерпретер, систем за управљање
	зависностима и тестови. Типичне акције таквих агената могу се груписати у
	неколико категорија:
	\begin{itemize}
		\item \textbf{навигација и претрага} — листање директоријума, претрага
		датотека по имену и садржају (нпр. \texttt{find}, \texttt{grep});
		\item \textbf{читање} — приказ садржаја датотеке или њеног дела;
		\item \textbf{измена} — креирање датотека и измена постојећег кода, нпр.
		заменом тачно одређеног дела текста;
		\item \textbf{извршавање} — покретање произвољних команди, скрипти за
		репродукцију грешке и тестова.
	\end{itemize}
	Произвољна команда љуске је у принципу довољна за све наведене акције,
	па се поставља питање да ли специјализовани алати доносе стварну предност
	или само повећавају контекст. Управо ово питање је предмет хипотеза овог
	рада.

	\subsection{Интерфејс између агента и рачунара}

	\textcite{yang2024sweagent} уводе појам интерфејса између агента и
	рачунара (ACI), по аналогији са корисничким интерфејсом за људе. Њихова
	теза је да језички модели имају другачије потребе од људи, па алати
	направљени за људе (нпр. интерактивни уређивачи текста) нису оптимални за
	агенте. Као принципе доброг интерфејса наводе једноставне акције са јасном
	документацијом, концизне и информативне повратне информације, као и
	заштитне механизме, нпр. проверу синтаксе након измене, који спречавају
	нагомилавање грешака. Аблацијом показују да дизајн интерфејса значајно
	утиче на успешност агента, иако је модел исти.

	\subsection{Изолација и поновљивост}

	Пошто агент извршава произвољне команде, окружење мора бити изоловано од
	рачунара на коме се систем покреће. Уобичајено решење су контејнери (нпр.
	Docker), који обезбеђују да сваки задатак почиње из истог, унапред
	дефинисаног стања, са инсталираним зависностима одговарајућих верзија.
	Изолација је истовремено безбедносни захтев и предуслов за поновљивост:
	разлике у окружењу не смеју да утичу на резултат поређења конфигурација.

	\section{Евалуациони оквир}

	\subsection{Бенчмаркови за код}

	Рани бенчмаркови, као што су HumanEval \parencite{chen2021humaneval} и
	MBPP \parencite{austin2021mbpp}, састоје се од кратких, самосталних
	функција које модел генерише на основу описа, а исправност се проверава
	тестовима. За такве задатке уведена је метрика pass@$k$, вероватноћа да је
	бар једно од $k$ генерисаних решења тачно. Ако се за задатак генерише
	$n \ge k$ решења, од којих је $c$ тачних, непристрасна процена износи
	\parencite{chen2021humaneval}:
	\begin{equation}
		\text{pass@}k = \mathbb{E}_{\text{задаци}}\left[1 - \frac{\binom{n-c}{k}}{\binom{n}{k}}\right].
	\end{equation}
	Такви задаци, међутим, не захтевају рад са постојећом кодном базом, што је
	суштина рада агената. Поред тога, \textcite{liu2023evalplus} показују да
	недовољни тестови доводе до прецењивања перформанси, јер нетачна решења
	пролазе проверу.

	\subsection{SWE-bench}

	SWE-bench \parencite{jimenez2024swebench} састоји се од 2294 задатка
	прикупљених из стварних проблема и одговарајућих захтева за спајање (енг.
	\textit{pull requests}) у 12 популарних Python репозиторијума. Сваки задатак
	садржи опис проблема и стање репозиторијума пре исправке. Систем треба да
	произведе исправку у облику разлике у односу на изворни код. Исправка се
	сматра успешном ако након њене примене:
	\begin{itemize}
		\item прођу тестови који су пре исправке падали, а које је увела
		оригинална исправка (FAIL\_TO\_PASS), и
		\item и даље пролазе тестови који су пролазили и пре исправке
		(PASS\_TO\_PASS), чиме се проверава да исправка није покварила постојећу
		функционалност.
	\end{itemize}
	Поред пуног скупа, често се користи SWE-bench Lite, подскуп од 300
	тест задатака који смањује трошак евалуације. Задаци су изабрани филтрирањем
	пуног скупа: уклоњени су задаци са сликама, спољашњим линковима и
	референцама на друге захтеве, задаци чији опис има мање од 40 речи, као и
	задаци чија референтна исправка мења више од једне датотеке, има више од
	три измењена блока (енг. \textit{hunks}) или креира и брише датотеке.
	Lite стога садржи самосталне, локализоване исправке грешака, што треба
	имати у виду при уопштавању резултата.

	Приликом тумачења резултата треба имати у виду познате проблеме скупа.
	\textcite{aleithan2024swebenchplus} показују да је у значајном делу
	успешних решења исправка већ садржана у опису проблема, као и да су
	поједини тестови превише слаби да открију нетачне исправке. Такође, пошто су
	репозиторијуми јавно доступни, постоји ризик да су задаци или њихова
	решења били део тренинг података модела (енг. \textit{data contamination}),
	што је мотивисало бенчмаркове попут LiveCodeBench
	\parencite{jain2024livecodebench}.

	\subsection{Метрике}

	Нека је $\mathcal{D}$ скуп од $N$ задатака, а $y_i \in \{0, 1\}$ индикатор
	да је задатак $i$ решен. Основна метрика успешности је стопа решених
	задатака (енг. \textit{resolve rate}):
	\begin{equation}
		R = \frac{1}{N} \sum_{i=1}^{N} y_i.
	\end{equation}

	Трошак задатка $i$ одређује се на основу броја улазних и излазних токена
	свих позива модела у трајекторији и цене по токену коју одређује провајдер:
	\begin{equation}
		C_i = N^{(i)}_{\text{in}} \cdot p_{\text{in}} + N^{(i)}_{\text{out}} \cdot p_{\text{out}},
	\end{equation}
	где $N_{\text{in}}$ следи из израза~\eqref{eq:ulazni_tokeni}. Код провајдера
	који подржавају кеширање, улазни токени се додатно деле на кеширане и
	некеширане, са различитим ценама. Из трошка се изводе просечан трошак по
	задатку $\bar{C} = \frac{1}{N}\sum_i C_i$ и трошак по решеном задатку:
	\begin{equation}
		C_{\text{решен}} = \frac{\sum_{i=1}^{N} C_i}{\sum_{i=1}^{N} y_i},
	\end{equation}
	која обухвата и трошак неуспешних покушаја. Допунске метрике су број корака
	(позива модела) по задатку, укупан број токена и време извршавања.

	Стопа решених задатака не показује да ли је агент бар пронашао место
	грешке. Зато се, по узору на \textcite{xia2025agentless}, посматра и
	локализација на нивоу датотеке. Нека је $F_i^{\text{агент}}$ скуп датотека
	које мења исправка агента, а $F_i^{\text{реф}}$ скуп датотека које мења
	референтна исправка задатка $i$. Стопа локализације је
	\begin{equation}
		L = \frac{1}{N} \sum_{i=1}^{N} \mathbb{1}\left[F_i^{\text{агент}} \cap F_i^{\text{реф}} \neq \emptyset\right].
	\end{equation}
	Пошто у скупу SWE-bench Lite референтна исправка увек мења тачно једну
	датотеку, $F_i^{\text{реф}}$ има један елемент, па метрика мери да ли је
	агент изменио управо ту датотеку. Ова метрика даје информацију и о
	задацима које агент није успео да реши,
	што је посебно корисно код слабијих модела, код којих је стопа решених
	задатака ниска.

	Пошто се конфигурације разликују и по успешности и по трошку, поређење се
	врши у дводимензионалном простору. Конфигурација $A$ доминира
	конфигурацијом $B$ ако није ни скупља ни мање успешна, а бар по једном
	критеријуму је строго боља. Скуп недоминираних конфигурација чини
	Парето фронт, који приказује најбоље доступне компромисе између цене и
	перформанси.

	\subsection{Статистичко поређење конфигурација}

	Резултати агената подложни су случајности, како због стохастичности модела,
	тако и због коначног броја задатака. \textcite{reimers2017distributions}
	показују да поређење заснивано на једном покретању може довести до
	погрешних закључака, па је пожељно извештавати расподелу резултата кроз
	више покретања. \textcite{card2020littlepower} указују на то да су многи
	експерименти у обради природног језика статистички подмоћни, тј. да број
	примера није довољан да се поуздано открије разлика очекиване величине.

	Када се две конфигурације евалуирају на истим задацима, опажања су
	упарена, па треба користити упарене тестове \parencite{dror2018hitchhiker}.
	За бинарни исход (решен/нерешен) одговарајући је Мекнемаров тест
	\parencite{mcnemar1947}, који разматра само неслагања: нека је $b$ број
	задатака које решава само конфигурација $A$, а $c$ број задатака које
	решава само конфигурација $B$. Под нултом хипотезом да су конфигурације
	подједнако успешне, $b$ има биномну расподелу $\mathrm{Bin}(b + c, 0{,}5)$,
	а за довољно велико $b + c$ користи се апроксимација
	\begin{equation}
		\chi^2 = \frac{(b - c)^2}{b + c}.
	\end{equation}
	За непрекидне метрике, као што је трошак по задатку, могу се користити
	Вилкоксонов тест ранга са знаком или пермутациони тестови. Интервали
	поверења за стопу решених задатака и разлике између конфигурација могу
	се проценити бутстреп методом \parencite{efron1993bootstrap}, поновљеним
	узорковањем задатака са враћањем. Исти поступак се примењује и на трошак по
	решеном задатку, који је количник збирова и нема вредност по појединачном
	задатку, па се за њега не може директно применити упарени тест.

	Неуспех статистичког теста да открије разлику не значи да разлика не
	постоји. Када је циљ показати да нова конфигурација није битно лошија од
	референтне, користи се тест неинфериорности: унапред се одређује маргина
	$\delta$, највећа разлика која се сматра прихватљивом, а неинфериорност се
	закључује ако доња граница интервала поверења за разлику $R_A - R_B$ лежи
	изнад $-\delta$. На тај начин се тврдња „конфигурације су подједнако
	успешне“ замењује проверљивом тврдњом „конфигурација $A$ није лошија од
	конфигурације $B$ за више од $\delta$“.
	
	% =========================== 4. МЕТОДОЛОГИЈА =================================
	\chapter{Методологија}
	
	Ово поглавље описује систем развијен за потребе рада и поступак којим је
	спроведен експеримент. Најпре је приказана архитектура агента, затим алати
	који му могу бити доступни и начин на који се бирају, окружење у коме се
	команде извршавају и, на крају, поступак евалуације: скуп података, дизајн
	експеримента, ток једног покретања, мере, правила искључивања и план
	статистичке анализе. Све вредности наведене у овом поглављу исте су у свим
	условима експеримента, осим скупа доступних алата.

	\section{Архитектура агента}

	Агент је реализован као петља позивања алата у духу парадигме ReAct
	\parencite{yao2023react}. Компоненте система и ток података између њих
	приказани су на Слици~\ref{fig:arhitektura}.

	\begin{figure}[H]
		\centering
		\includegraphics[width=\textwidth]{dijagrami/arhitektura.png}
		\caption{Архитектура развијеног агентског система}
		\label{fig:arhitektura}
	\end{figure}

	Један корак петље састоји се од следећих радњи:
	\begin{enumerate}
		\item Модел добија историју поруке (системски упит, опис задатка и све
		досадашње позиве алата са резултатима) и JSON шеме активних алата.
		\item Ако одговор модела не садржи позив алата, покретање се завршава
		(разлог завршетка \texttt{completed}).
		\item У супротном се аргументи сваког позива валидирају, позив се
		извршава у изолованом окружењу, а резултат се враћа моделу као порука
		улоге \texttt{tool}.
		\item Петља се понавља највише \texttt{max\_turns} корака, након чега се
		покретање прекида (разлог завршетка \texttt{max\_turns}).
	\end{enumerate}

	Модел ништа не извршава директно, већ само предлаже позив алата, а систем
	одлучује да ли ће га и како извршити. Грешке су део разговора: неисправни
	аргументи, непознат алат или неуспешна команда не прекидају петљу, већ се
	моделу враћају као текстуална порука облика \texttt{ERROR: ...}, на коју
	модел може да реагује. Излаз алата се тиме третира као део упита, па је
	јасноћа порука о грешкама део дизајна система. Ток једног покретања кроз
	петљу, укључујући обраду грешака, приказан је на Слици~\ref{fig:petlja}.

	\begin{figure}[H]
		\centering
		\includegraphics[height=0.8\textheight]{dijagrami/petlja.png}
		\caption{Петља агента: позиви модела, извршавање алата и услови завршетка}
		\label{fig:petlja}
	\end{figure}

	Комуникација са моделом одвојена је од петље апстракцијом провајдера, која
	подржава OpenAI и Azure OpenAI API. Петља зна само за провајдерски
	неутралан облик одговора (текст, листа позива алата и потрошња токена), па
	се исти агент може користити са различитим провајдерима без измене логике.

	Систем током рада бележи телеметрију: број улазних, излазних и токена
	закључивања, цену израчунату на основу ценовника модела, број позива
	модела, време проведено у позивима модела и у извршавању алата, као и број
	позива и грешака по алату. Провајдер повремено одбија исправан захтев због
	лажно позитивне класификације садржаја (\texttt{invalid\_prompt}); такав
	захтев се понавља највише два пута, а свака друга грешка API-ја прекида
	покретање. Петља подржава и кооперативни сигнал за заустављање, који
	користи спољашње ограничење укупног трајања покретања.

	Контекст се не сажима: историја поруке расте током целог покретања. Као
	што је показано у изразу~\eqref{eq:ulazni_tokeni}, број улазних токена
	тада расте приближно квадратно са бројем корака, па је ограничавање
	величине излаза алата важан део контроле трошка.

	\section{Дизајн алата}

	Агент у експерименту има на располагању највише пет алата. Модел сваки алат
	види као опис функције са шемом параметара, а та шема се шаље уз сваки позив
	модела, без обзира на то да ли ће алат у том кораку бити употребљен. Пре
	извршавања систем проверава аргументе према шеми. Неисправни аргументи или
	непознат алат не прекидају рад: грешка се враћа моделу као обичан
	текстуални резултат, на који он може да реагује у следећем кораку.

	Сви алати раде у истом изолованом контејнеру, у директоријуму \texttt{/work},
	у коме се налази изворни код пројекта. Алати \texttt{bash} и \texttt{search}
	своде се на једну команду љуске, а \texttt{str\_replace}, \texttt{read\_file}
	и \texttt{find\_file} у контејнеру покрећу кратку скрипту. Излаз сваког
	алата ограничен је на 8000 знакова, а дужи излаз се сече и на крај се додаје
	напомена да је приказ непотпун. Преглед алата дат је у
	Табели~\ref{tab:alati}. Пуне JSON шеме свих алата, у облику у коме се шаљу
	моделу, доступне су у изворном коду система, у директоријуму
	\texttt{src/harness/tools} на грани \texttt{main}
	репозиторијума \url{https://github.com/aleksaljujic/agent-harness}.

	\begin{table}[H]
		\caption{Алати агента}
		\label{tab:alati}
		\tabelalevo
		\begin{tabular}{|l|>{\raggedright\arraybackslash}p{2.7cm}|>{\raggedright\arraybackslash}p{2.5cm}|>{\raggedright\arraybackslash}p{2.3cm}|>{\raggedright\arraybackslash}p{1.9cm}|l|}
			\hline
			Алат & Намена & Обавезни параметри & Опциони параметри & Ограничење излаза & Шема \\
			\hline
			\texttt{bash} & произвољна команда љуске & команда, образложење & --- & 8000 знакова, 60~s & 37 \\
			\hline
			\texttt{str\_replace} & измена постојеће датотеке & путања, стари текст, нови текст & све појаве, редни број појаве & --- & 223 \\
			\hline
			\texttt{search} & претрага садржаја по регуларном изразу & образац & филтер датотека & 50 погодака & 43 \\
			\hline
			\texttt{read\_file} & читање датотеке са бројевима линија & путања & почетна и завршна линија & 8000 знакова & 91 \\
			\hline
			\texttt{find\_file} & претрага датотека по имену & образац & највећи број резултата & 100 резултата & 105 \\
			\hline
		\end{tabular}
	\end{table}

	Колона „Шема“ даје број улазних токена које опис алата додаје сваком позиву
	модела. Вредности су измерене на моделу \texttt{gpt-5.4-nano}, као разлика
	броја улазних токена истог захтева послатог са алатом и без њега. Кад је у
	листи бар један алат, томе се једном по позиву додаје још 98 токена за оквир
	листе. Услови из нацрта $2 \times 2$ тако носе по позиву: А 135, Б 358, В 374
	и Г 597 токена. Сам \texttt{str\_replace} (223) кошта скоро колико цела група
	\texttt{search}, \texttt{read\_file} и \texttt{find\_file} заједно (239), јер
	његов опис садржи упутства за случај вишеструког поклапања. Систем има и
	шести алат, \texttt{run\_tests}, за покретање тестова са структурисаним
	резултатом, али он није био део експеримента, јер у овој поставци агент ради
	без покретања тестова пројекта.

	\subsection{Алат bash}

	Алат \texttt{bash} је основни алат и једини који је присутан у сваком
	услову. Сам по себи довољан је за решавање задатка: кроз њега агент може да
	претражује, чита и мења датотеке, али без икакве контроле над обликом и
	обимом излаза. Алат прима команду љуске и враћа спојен стандардни излаз и
	излаз за грешке.

	Шема има два обавезна параметра: саму команду и поље за образложење. У то
	поље модел пре сваког позива кратко наводи зашто покреће команду. Систем
	образложење не користи при извршавању, али оно остаје у историји разговора,
	па га модел види у наредним корацима, а сачувано је и у транскриптима
	покретања.

	Свака команда извршава се као нов процес љуске у \texttt{/work}. Стање попут
	промене директоријума или променљивих окружења зато се не преноси из једног
	позива у други, па модел зависне кораке мора да споји у једну команду.

	Излаз дужи од 8000 знакова се сече, а на крај се додаје напомена о томе
	колико је знакова изостављено, уз савет да се команда сузи, како модел
	непотпун приказ не би схватио као целину. Команда која траје дуже од 60
	секунди се прекида, а модел добија поруку о истеку времена. Ако команда не
	произведе никакав излаз, враћа се само њен излазни код. Тако се празан
	резултат разликује од грешке, а модел зна да ли је тиха команда успела.

	\subsection{Алат str\_replace}

	Алат \texttt{str\_replace} мења постојећу датотеку тако што тачно наведени
	текст замени новим, без поновног генерисања целе датотеке. Обавезни
	параметри су путања, стари и нови текст. Поклапање је дословно, а не по
	регуларном изразу, па знакови који су у коду чести, попут тачке, звездице и
	заграда, не морају да се означавају.

	Подразумевано, стари текст мора да се појави тачно једном. Ако га нема,
	модел добија грешку са саветом да прочита датотеку и препише текст заједно
	са увлачењем. Ако се појављује више пута, измена се не врши, а порука наводи
	број појава и број линије сваке од њих. Порука уз то саветује да се стари
	текст прошири навише, до претходног коментара или дефиниције функције или
	класе, јер су поновљени блокови кода обично једнаки и испод, па ширење
	наниже ретко разрешава двосмисленост. Ова провера је безбедносно својство
	алата, јер би без ње измена тихо погодила погрешно место.

	За намерно вишеструке измене постоје два опциона параметра. Један мења све
	појаве, а други само одређену појаву по редоследу у датотеци, бројано од
	прве. Параметри се међусобно искључују, а њихова комбинација, као и редни
	број ван опсега, враћа грешку са стварним бројем појава. После успешне
	измене модел добија број замењених појава и линије на којима су биле.

	Алат постоји зато што је ово једина операција коју \texttt{bash} ради лоше.
	Програм \texttt{sed} ради ред по ред, па не може да обухвати блок од више
	линија, а сваки посебан знак из кода мора да се означи. Над обрасцем који се
	јавља више пута \texttt{sed} тихо мења све појаве, без упозорења.

	\subsection{Алат search}

	Алат \texttt{search} претражује садржај датотека у целом пројекту по
	проширеном регуларном изразу. Обавезан параметар је образац. Опциони филтер
	има два значења: ако не садржи косу црту, тумачи се као образац имена
	датотеке, на пример све Python датотеке, а ако садржи косу црту, тумачи се
	као путања датотеке или директоријума на коју се претрага ограничава. Ова
	подела постоји зато што се филтер по имену пореди само са основним именом
	датотеке, па би путања као филтер тихо промашила све датотеке.

	Сваки погодак враћа се у једном реду: путања, број линије и садржај линије.
	Тај облик непосредно је употребљив за \texttt{read\_file} и
	\texttt{str\_replace}. Директоријуми са историјом верзија (\texttt{.git}) и
	инсталираним зависностима (\texttt{node\_modules}) се прескачу, јер у њима
	нема кода који агент треба да мења, а стварају много шума.

	Алат враћа највише 50 погодака, редоследом обиласка система датотека, а не
	по релевантности. Када погодака има више, остали се одбацују без икакве
	назнаке: модел не добија ни укупан број погодака ни ознаку да је листа
	скраћена. По томе се \texttt{search} разликује од \texttt{find\_file}. Уз то
	важи и општа граница од 8000 знакова, па дугачке линије могу да пресеку
	излаз и пре педесетог погодка, овај пут уз напомену. Када нема погодака,
	модел добија изричиту поруку уместо празног излаза, а неисправан регуларни
	израз враћа поруку о грешци.

	Алат не даје агенту ништа што \texttt{bash} већ не може. Његова вредност је у
	ограничавању излаза, јер необуздана рекурзивна претрага над правим
	репозиторијумом може једним позивом да попуни контекст.

	\subsection{Алат read\_file}

	Алат \texttt{read\_file} чита датотеку и враћа њен садржај са бројевима
	линија. Обавезан параметар је путања, а опциони су почетна и завршна линија,
	које се броје од прве и улазе у опсег. Без њих се враћа цела датотека.

	Сваки ред излаза почиње бројем линије и табулатором. Бројеви линија повезују
	алате: погоци из \texttt{search} и грешке из \texttt{str\_replace} наводе
	линије, па агент може одмах да затражи тачан опсег. Број линије није део
	датотеке, па опис алата \texttt{str\_replace} изричито упозорава да се тај
	префикс уклони пре него што се текст препише у стари текст, јер у супротном
	тражени текст неће бити пронађен.

	Опсег се проверава пре читања. Грешку враћају почетна линија мања од прве,
	почетна линија после завршне и почетна линија после краја датотеке, при чему
	порука у последњем случају наводи и колико линија датотека има. Завршна
	линија после краја датотеке није грешка, него се тихо своди на последњу
	линију, јер је читање до краја уобичајена намера. У томе је разлика у односу
	на \texttt{bash}, где исечак датотеке са погрешним опсегом не враћа ништа
	осим излазног кода, док овде модел добија грешку по којој може да поступи.

	Посебне поруке постоје и за непостојећу датотеку, за директоријум и за
	датотеку која није у кодирању UTF-8, а празна датотека даје изричиту поруку
	да је празна. Алат нема сопствено ограничење броја линија, па важи општа
	граница од 8000 знакова. Велика датотека прочитана без опсега се сече уз
	напомену, и агент тада треба да је чита по деловима.

	\subsection{Алат find\_file}

	Алат \texttt{find\_file} тражи датотеке по имену или путањи, а не по
	садржају, за шта служи \texttt{search}. Обавезан параметар је образац, а
	опциони параметар одређује највећи број резултата и подразумевано износи
	100.

	Ако образац садржи џокер знакове (звездицу, упитник или угласту заграду),
	тумачи се као глоб и пореди се и са именом датотеке и са њеном путањом у
	односу на корен пројекта. Тако раде и обрасци по имену, на пример све Python
	датотеке, и обрасци по путањи, на пример тестови у одређеном
	поддиректоријуму. Образац без џокер знакова тумачи се као део имена
	датотеке, уз разликовање великих и малих слова. Враћају се само датотеке, не
	и директоријуми, а као и код \texttt{search}, прескачу се \texttt{.git} и
	\texttt{node\_modules}.

	Резултат је списак релативних путања, по једна у реду, абецедно уређен, па
	исти упит увек даје исти редослед. Када погодака има више од граничне
	вредности, приказује се првих N, а на крају се додаје ред са укупним бројем
	погодака, у облику „приказано првих N од M“. За разлику од \texttt{search},
	модел тако зна да је списак непотпун и колико треба да сузи образац. Када
	погодака нема, модел добија изричиту поруку.

	Као и \texttt{search}, овај алат не доноси нову могућност, јер исто може да
	се уради кроз \texttt{bash}. Разлика је у облику излаза: програм
	\texttt{find} има заставице које модел мора сваки пут да погоди, а излаз му
	није ограничен, док \texttt{find\_file} увек даје исти, ограничен и уређен
	облик резултата.

	\section{Окружење за извршавање}

	Свако покретање добија сопствени, једнократни Docker контејнер.
	Коришћена слика садржи Python 3.12 и основне алате, али не и зависности
	пројекта над којим агент ради. Агент зато не може да покрене тестове
	пројекта, већ исправку прави на основу читања кода и описа проблема. Из
	истог разлога алат \texttt{run\_tests} није део експеримента.
	
	Јединствена слика за све репозиторијуме изабрана је због упоредивости
	услова и изводљивости експеримента: окружење је исто за сваки задатак, па
	разлике у зависностима не могу да утичу на поређење услова, а 1200
	покретања остаје у оквиру расположивог времена и средстава. Припрема
	засебног окружења по задатку захтевала би вишеструко више простора и
	времена.
	
	Последица те одлуке је да агент не може да покрене тестове пројекта нити да
	репродукује грешку, већ ради на основу читања кода и описа проблема. Успут
	се уклања и могућност да агент подешава исправку према постојећим
	тестовима, мада главну заштиту од доласка до званичног решења чине
	искључен приступ мрежи и уклоњена историја верзија, пошто скривени тестови
	којима се исправка оцењује ионако нису део радне копије.

	Контејнер се покреће са ограниченим животним веком, тако да се уклања и ако
	систем неочекивано прекине рад, а команде прима преко \texttt{docker exec}.
	Систем датотека је трајан између позива, али стање љуске није: свака
	команда почиње у новој љусци, што је наведено и у системском упиту.
	Радна копија репозиторијума повезана је са контејнером као директоријум
	\texttt{/work} (енг. \textit{bind mount}), па су измене агента одмах видљиве
	на рачунару домаћину, одакле се издваја исправка.

	Контејнер ради са следећим ограничењима:
	\begin{itemize}
		\item без приступа мрежи (\texttt{--network none}), па агент не може да
		преузме званичну исправку проблема;
		\item највише 2~GB меморије и 256 процеса;
		\item процеси раде под корисником домаћина, а не као \texttt{root};
		\item скрипте алата доступне су само за читање;
		\item излаз команде ограничен је на 8000 знакова, уз изричиту ознаку
		скраћивања, а подразумевано време извршавања једне команде на 60~s.
	\end{itemize}
	На нивоу покретања, агент има највише 40 корака и 1800~s укупног времена.

	\section{Поступак евалуације}

	\subsection{Пилот евалуација}

	Током развоја систем је провераван на 12 малих, синтетички направљених
	задатака (креирање и измена датотека, исправљање грешака, рефакторисање,
	рад са тестовима и CSV датотекама). Сваки задатак се извршава у новом
	контејнеру, а успех се утврђује провером стања система датотека након
	покретања, а не на основу онога што модел тврди да је урадио. Пошто су ови
	задаци коришћени за развој и подешавање система, а разликују се од стварних
	проблема из репозиторијума и по величини и по сложености, резултати пилота
	служе само као мотивација за главни експеримент и не представљају доказ за
	хипотезе.

	\subsection{Скуп података}

	Главни експеримент користи скуп SWE-bench Lite
	\parencite{jimenez2024swebench}, у \texttt{test} подели, описан у поглављу~3, и то свих 300 задатака из 12 репозиторијума. Највише
	задатака потиче из репозиторијума \texttt{django} (114) и \texttt{sympy}
	(77). Пошто се користи цео скуп, нема избора задатака од стране аутора, па
	ни пристрасности при избору. Сваки задатак обухвата репозиторијум, ознаку
	комита пре исправке (\texttt{base\_commit}) и текст проблема. Скривени
	тестови (FAIL\_TO\_PASS и PASS\_TO\_PASS) агенту нису доступни.

	\subsection{Дизајн експеримента}

	Експеримент има факторски дизајн $2 \times 2$ са два фактора: присуство
	алата \texttt{str\_replace} и присуство групе алата за претрагу и читање
	(\texttt{search}, \texttt{read\_file} и \texttt{find\_file}). Алат
	\texttt{bash} доступан је у свим условима. Комбинацијом нивоа добијају се
	четири услова приказана у Табели~\ref{tab:uslovi}.

	\begin{table}[H]
		\caption{Услови експеримента}
		\label{tab:uslovi}
		\tabelalevo
		\begin{tabular}{|l|l|l|}
			\hline
			& без \texttt{str\_replace} & са \texttt{str\_replace} \\
			\hline
			без претраге и читања & А: \texttt{bash} & Б: \texttt{bash}, \texttt{str\_replace} \\
			\hline
			са претрагом и читањем & В: \texttt{bash}, претрага & Г: сви алати \\
			\hline
		\end{tabular}
	\end{table}

	Сваки задатак пролази кроз сва четири услова, по једном, што даје 1200
	покретања. Дизајн је упарен: задатак представља блок, па се услови пореде
	на истим задацима. Ефекат алата \texttt{str\_replace} добија се поређењем
	услова А и Б, као и В и Г, а ефекат групе алата за претрагу и читање
	поређењем услова А и В, као и Б и Г.

	Све остале променљиве су фиксиране и приказане у
	Табели~\ref{tab:parametri}. Посебно, сви услови користе исти системски
	упит, који не помиње ниједан конкретан алат нити обесхрабрује одређени
	начин измене датотека. На тај начин се мери ефекат доступности алата, а не
	ефекат алата у комбинацији са упутством да се он користи.

	\begin{table}[H]
		\caption{Фиксирани параметри експеримента}
		\label{tab:parametri}
		\tabelalevo
		\begin{tabular}{|l|l|}
			\hline
			Параметар & Вредност \\
			\hline
			модел & \texttt{gpt-5.4-nano} \\
			\hline
			температура & 0 \\
			\hline
			режим проширеног закључивања & искључен \\
			\hline
			највећи број корака & 40 \\
			\hline
			највеће трајање покретања & 1800~s \\
			\hline
			приступ мрежи & искључен \\
			\hline
			број понављања по услову & 1 \\
			\hline
			системски упит & идентичан, без помињања алата \\
			\hline
		\end{tabular}
	\end{table}

	\subsection{Ток једног покретања}

	Једно покретање састоји се од четири фазе:
	\begin{enumerate}
		\item \textbf{Припрема.} Репозиторијум се из локалне копије поставља на
		стање комита пре исправке (\texttt{base\_commit}). Историја верзија се затим брише и замењује
		једним почетним комитом без удаљеног репозиторијума, тако да агент не
		може да дође до касније, званичне исправке.
		\item \textbf{Рад агента.} Агент добија текст проблема као задатак и
		ради над радном копијом у контејнеру.
		\item \textbf{Издвајање исправке.} Након завршетка, исправка се добија као
		разлика радне копије у односу на почетни комит. Бележе се и
		дијагностички подаци (изостала исправка, број измењених датотека и линија),
		а исправка се одмах уписује у званичном SWE-bench формату предвиђања.
		\item \textbf{Оцењивање.} У одвојеној фази, званична SWE-bench процедура
		примењује исправку у Docker слици припремљеној за тај задатак и покреће
		скривене тестове. Задатак је решен ако пролазе сви FAIL\_TO\_PASS и
		PASS\_TO\_PASS тестови. Слике окружења за оцењивање, приближно 4~GB по
		задатку, преузимају се унапред, јер евалуатор одустаје од преузимања
		после десет минута и задатак тада добија лажно негативну оцену. Задаци
		код којих агент није произвео исправку не прослеђују се евалуатору и
		рачунају се као нерешени. Ако оцена не може да се добије из других
		разлога, то се бележи засебно и не третира као неуспех агента.
	\end{enumerate}

	Ток експеримента у целини, од припреме репозиторијума до израчунавања мера,
	приказан је на Слици~\ref{fig:tok-eksperimenta}. Исти задатак пролази кроз
	сва четири услова, а сваки услов ради над засебном радном копијом.

	\begin{figure}[H]
		\centering
		\includegraphics[width=0.7\textwidth]{dijagrami/tok_eksperimenta.png}
		\caption{Ток експеримента}
		\label{fig:tok-eksperimenta}
	\end{figure}

	За свако покретање чувају се ред са метрикама, цео разговор са моделом у
	читљивом облику, необрађене поруке које омогућавају тачно понављање
	анализе и детаљи евентуалне грешке. Уз сваки ред бележи се и ознака
	верзије кода система (\texttt{harness\_sha}). Експеримент је започет пре
	него што су последње измене кода биле комитоване, па ова ознака указује на
	претходни комит. Измене које мењају понашање агента унете су пре почетка
	експеримента, а измене начињене током њега односе се искључиво на
	чување резултата и наставак прекинутог рада, тако да су сва покретања
	међусобно упоредива.

	\subsection{Мере}

	Мере су унапред подељене на три групе, како би се избегло накнадно
	тражење статистички значајних разлика међу великим бројем колона:
	\begin{itemize}
		\item \textbf{Примарне:} стопа решених задатака, трошак по решеном
		задатку (укупан трошак услова подељен бројем решених задатака у том
		услову) и стопа локализације, дефинисане у поглављу~3;
		\item \textbf{Секундарне:} трошак по покретању, број улазних и излазних
		токена, број корака, односно позива модела, број позива алата, укупно време, време у позивима
		модела и време извршавања алата;
		\item \textbf{Дијагностичке:} изостала исправка, разлог завршетка покретања,
		грешке алата, број позива по алату, број FAIL\_TO\_PASS и PASS\_TO\_PASS
		тестова који пролазе и величина исправке.
	\end{itemize}

	Трошак се рачуна према изразу из поглавља~3, на основу броја токена и
	јавног ценовника модела \texttt{gpt-5.4-nano} (0,20 долара по милиону
	улазних и 1,25 долара по милиону излазних токена). Стопа локализације
	рачуна се накнадно, поређењем датотека које мења исправка агента са
	датотеком коју мења референтна исправка; изостала исправка се рачуна као
	неуспешна локализација, по узору на \textcite{xia2025agentless}. Број позива по алату показује да ли је доступан
	алат заиста коришћен, чиме се ефекат доступности алата разликује од
	ефекта његовог коришћења.

	\subsection{Правила искључивања}

	Правила искључивања утврђена су пре увида у резултате:
	\begin{itemize}
		\item одбијање захтева због класификације садржаја
		(\texttt{invalid\_prompt}) у примарној анализи сматра се исходом
		покретања, тј. нерешеним задатком, а стопа одбијања извештава се по
		услову;
		\item достизање највећег броја корака или највећег трајања такође је
		исход (нерешен задатак), а не квар;
		\item искључују се само покретања прекинута због грешке система или
		инфраструктуре. Јединица искључивања је задатак: ако је покретање
		искључено у једном услову, задатак се искључује из свих услова, како би
		број задатака по услову остао исти;
		\item анализа осетљивости понавља поређења само на задацима који су у
		сва четири услова завршени без грешке API-ја.
	\end{itemize}

	\subsection{Статистичка анализа}

	Хипотезе Х1 и Х2 проверавају се на главним ефектима факторског дизајна, а
	хипотеза Х3 поређењем услова Б и Г:
	\begin{itemize}
		\item \textbf{Х1 (трошак по решеном задатку).} Упоређује се трошак по
		решеном задатку услова са алатом \texttt{str\_replace} (Б и Г) и без
		њега (А и В). Пошто је ова мера количник збирова, 95\% интервал
		поверења за разлику процењује се бутстреп методом, поновљеним
		узорковањем задатака са враћањем, при чему се задатак узоркује
		истовремено у свим условима. Као допунски упарени тест, за сваки задатак
		рачуна се разлика просечног трошка по покретању са алатом и без њега и
		примењује Вилкоксонов тест ранга са знаком. Хипотеза се потврђује ако цео
		интервал поверења лежи испод нуле.
		\item \textbf{Х2 (трошак и локализација).} Разлика у трошку по покретању
		између услова са групом алата за претрагу и читање (В и Г) и без ње (А и
		Б) тестира се Вилкоксоновим тестом над разликама по задатку. Разлика у
		стопи локализације процењује се упареним бутстреп интервалом поверења,
		уз Мекнемаров тест за појединачна поређења А--В и Б--Г. Део Х2а
		потврђује се ако је медијана разлике позитивна, а $p$-вредност после
		Холмове корекције мања од 0,05, а део Х2б ако цео интервал поверења лежи
		изнад нуле. Хипотеза Х2 потврђена је само ако важе оба дела.
		\item \textbf{Х3 (неинфериорност).} За разлику стопа решених задатака
		$R_{\text{Б}} - R_{\text{Г}}$ рачуна се 95\% упарени бутстреп интервал
		поверења \parencite{efron1993bootstrap}. Неинфериорност се закључује ако
		је доња граница интервала изнад минус пет процентних поена.
	\end{itemize}

	Сви бутстреп интервали добијају се из 10\,000 понављања, као перцентилни
	интервали, са фиксним почетним стањем генератора случајних бројева, па су
	међусобно упоредиви. Ако у неком бутстреп узорку група нема ниједан решен
	задатак, трошак по решеном задатку није дефинисан и узорак се одбацује; ако
	је таквих узорака више од 5\%, хипотеза Х1 извештава се само описно.
	
	Примарну фамилију тестова, на коју се примењује Холмова корекција, чине
	четири теста: Вилкоксонови тестови за Х1 и Х2а и Мекнемарови тестови стопе
	локализације за парове А--В и Б--Г. За стопу решених задатака егзактни
	Мекнемаров тест \parencite{mcnemar1947} примењује се на сва четири контраста
	нацрта (А--Б, В--Г, А--В и Б--Г), у складу са препоруком да се за упарене
	податке користе упарени тестови \parencite{dror2018hitchhiker}. Ти тестови
	чине засебну, секундарну фамилију са сопственом Холмовом корекцијом и не
	служе за потврду хипотеза. Уз $p$-вредности се увек
	извештавају величине ефекта и интервали поверења. Интеракција фактора, као
	и разлике у ефекту по репозиторијуму, величини референтне исправке и броју
	FAIL\_TO\_PASS тестова, анализирају се експлоративно.
	
	% =================== 5. ОПИС И РЕЗУЛТАТИ ЕКСПЕРИМЕНТА ========================
	\chapter{Опис и резултати експеримента}
	
	Ово поглавље приказује поставку спроведеног експеримента и добијене
	резултате. Хипотезе Х1 и Х3 ослањају се на стопу решених задатака,
	утврђену званичним SWE-bench евалуатором, док се обе мере хипотезе Х2,
	трошак по покретању и стопа локализације, добијају без покретања тестова.
	Резултати су приказани редоследом хипотеза, а затим су дискутовани
	заједно.
	
	\section{Поставка експеримента}
	
	Поређена су четири услова добијена укрштањем два фактора: доступности алата
	\texttt{str\_replace} за измену датотека и доступности групе алата за
	претрагу (\texttt{search}, \texttt{read\_file} и \texttt{find\_file}). Алат
	\texttt{bash} доступан је у свим условима, а услови су приказани у
	Табели~\ref{tab:uslovi}.
	
	Експеримент обухвата цео \texttt{test} подскуп збирке SWE-bench Lite
	(300 задатака) у сва четири услова, дакле $300 \times 4 = 1200$ покретања,
	са једним покретањем по услову. Агент ради у контејнеру без приступа мрежи,
	са највише 40 корака и ограничењем од 1800 секунди по задатку. Укупно
	утрошено време износи 15,3 сата, а укупан обрачунати трошак \$33,13, о чему
	се детаљније говори међу ограничењима рада.
	
	Оцењивање решености спроведено је након завршетка свих покретања, на
	засебном рачунару са 64~GB радне меморије, без поновног покретања агента и
	без додатних позива модела. Слике окружења за све задатке, укупно
	приближно 100~GB, преузете су пре почетка оцењивања, а само оцењивање
	спроведено је са осам паралелних процеса. У услову Г евалуатор за два
	задатка није дао резултат тестова, па се они рачунају као нерешени, што је
	размотрено међу ограничењима рада.
	
	\section{Резултати}
	
	\subsection{Провера хипотезе Х1}
	
	Хипотеза Х1 тврди да доступност алата \texttt{str\_replace} смањује трошак
	по решеном задатку. Трошак по решеном задатку дефинисан је као однос укупног
	трошка свих покретања у групи и броја решених задатака у тој групи, чиме се
	у меру урачунава и трошак неуспешних покретања. Поређен је главни ефекат
	фактора, група услова са алатом (Б и Г) наспрам групе услова без њега (А и
	В), при чему свака група садржи по један услов са алатима за претрагу и
	један без њих. Стопе решених задатака и трошкови по условима приказани су у
	Табели~\ref{tab:h1_uslovi}, а трошак по решеном задатку са интервалима
	поверења и на Слици~\ref{fig:h1-uslovi}.
	
	\begin{table}[H]
		\caption{Трошак по решеном задатку по условима (300 задатака по услову)}
		\label{tab:h1_uslovi}
		\tabelalevo
		\begin{tabular}{|l|l|l|l|l|l|}
			\hline
			Услов & Алати & Решено & Трошак & Трошак/покр. & Трошак/решен \\
			\hline
			А & \texttt{bash} & 111 (37,0\%) & \$6,27 & \$0,0209 & \$0,0565 \\
			\hline
			Б & + \texttt{str\_replace} & 119 (39,7\%) & \$6,97 & \$0,0232 & \$0,0585 \\
			\hline
			В & + претрага & 105 (35,0\%) & \$10,56 & \$0,0352 & \$0,1006 \\
			\hline
			Г & сви алати & 84 (28,0\%) & \$9,34 & \$0,0311 & \$0,1111 \\
			\hline
		\end{tabular}
	\end{table}
	
	\begin{figure}[H]
		\centering
		\includegraphics[width=0.95\textwidth]{slike/h1_uslovi.png}
		\caption{Трошак по решеном задатку по условима, са 95\% бутстреп интервалима поверења}
		\label{fig:h1-uslovi}
	\end{figure}
	
	\textbf{Х1 --- одбачена.} Трошак по решеном задатку у групи са алатом
	\texttt{str\_replace} износи \$0,0803 (\$16,30 за 203 решена задатка), а у
	групи без њега \$0,0779 (\$16,83 за 216 решених). Разлика износи
	$+\$0{,}0024$ ($+3{,}1\%$), са упареним бутстреп интервалом поверења од
	95\% $[-\$0{,}0058,\ +\$0{,}0110]$ (10\,000 понављања, поновно узорковање
	задатака; ниједан узорак није одбачен због групе без решених задатака).
	Критеријум потврде био је да цео интервал лежи испод нуле, а добијени
	интервал обухвата нулу, при чему тачкаста процена има супротан смер од
	очекиваног. Интервали поверења за примарну анализу и анализу осетљивости
	приказани су на Слици~\ref{fig:h1-interval}, а бутстреп расподела разлике
	на Слици~\ref{fig:h1-bootstrap}.
	
	\begin{figure}[H]
		\centering
		\includegraphics[width=0.95\textwidth]{slike/h1_interval.png}
		\caption{Интервали поверења за разлику трошка по решеном задатку између услова са алатом \texttt{str\_replace} и услова без њега}
		\label{fig:h1-interval}
	\end{figure}
	
	\begin{figure}[H]
		\centering
		\includegraphics[width=0.95\textwidth]{slike/h1_bootstrap.png}
		\caption{Бутстреп расподела разлике трошка по решеном задатку у примарној анализи}
		\label{fig:h1-bootstrap}
	\end{figure}
	
	Допунски Вилкоксонов тест упарених рангова, спроведен над разликом трошка
	по покретању, даје медијану разлике $-\$0{,}00103$, $W = 19\,799$,
	$p = 0{,}065$ и $r_{rb} = -0{,}12$, и не опстаје после Холмове корекције
	примарне фамилије. Покретања са алатом \texttt{str\_replace} су, дакле, у
	просеку нешто јефтинија, за 3,1\%, али је и стопа решених задатака у тој
	групи нижа (33,8\% наспрам 36,0\%), па се уштеда по покретању поништава на
	нивоу решеног задатка. Ово разлагање приказано је на
	Слици~\ref{fig:h1-razlaganje}.
	
	\begin{figure}[H]
		\centering
		\includegraphics[width=0.95\textwidth]{slike/h1_razlaganje.png}
		\caption{Разлагање трошка по решеном задатку на трошак по покретању и стопу решених задатака}
		\label{fig:h1-razlaganje}
	\end{figure}
	
	Анализа осетљивости, ограничена на 256 задатака на којима ни у једном
	услову није дошло до одбијања захтева од стране провајдера, даје исти
	закључак: разлика износи $+\$0{,}0002$ ($+0{,}3\%$), са интервалом
	$[-\$0{,}0066,\ +\$0{,}0076]$. У овом подскупу Вилкоксонов тест јесте
	значајан ($p = 0{,}015$, $r_{rb} = -0{,}18$) и опстаје после Холмове
	корекције, али се односи на трошак по покретању: покретања са алатом
	јефтинија су за 5,9\%, а стопа решених нижа за 6,2\%, тако да трошак по
	решеном задатку остаје непромењен. Неслагање примарне анализе и анализе
	осетљивости у Холмовом поступку пријављује се као такво, а на закључак о
	хипотези Х1 не утиче, пошто је критеријум потврде интервал поверења.
	
	\subsection{Мера локализације}
	
	Као примарна мера за хипотезу Х2 користи се локализација, која не зависи од
	покретања тестова: покретање се сматра
	локализованим ако скуп датотека које мења исправка агента и скуп датотека које
	мења референтна исправка имају непразан пресек. Изостала исправка се по дефиницији
	сматра нелокализованом. Мера је укључена у план анализе пре првог
	израчунавања.
	
	У Табели~\ref{tab:rezultati} приказани су резултати по условима, а у
	Табели~\ref{tab:ishodi} исходи покретања.
	
	\begin{table}[H]
		\caption{Резултати по условима (300 задатака по услову)}
		\label{tab:rezultati}
		\tabelalevo
		\begin{tabular}{|l|l|l|l|l|l|}
			\hline
			Услов & Алати & Локализација & Трошак/зад. & Корака & Токена \\
			\hline
			А & \texttt{bash} & 78,0\% & \$0,0209 & 16,0 & 96\,550 \\
			\hline
			Б & + \texttt{str\_replace} & 77,3\% & \$0,0232 & 15,8 & 108\,896 \\
			\hline
			В & + претрага & 74,3\% & \$0,0352 & 23,9 & 168\,070 \\
			\hline
			Г & сви алати & 71,7\% & \$0,0311 & 22,4 & 149\,454 \\
			\hline
		\end{tabular}
	\end{table}
	
	\begin{table}[H]
		\caption{Исходи покретања по условима}
		\label{tab:ishodi}
		\tabelalevo
		\begin{tabular}{|l|l|l|l|l|}
			\hline
			Услов & Алати & Изостала & \texttt{invalid\_prompt} & Исцрпљени \\
			 &  & исправка &  & кораци \\
			\hline
			А & \texttt{bash} & 22 & 23 & 4 \\
			\hline
			Б & + \texttt{str\_replace} & 31 & 29 & 6 \\
			\hline
			В & + претрага & 38 & 6 & 41 \\
			\hline
			Г & сви алати & 52 & 3 & 54 \\
			\hline
		\end{tabular}
	\end{table}
	
	Однос између просечног трошка по задатку и стопе локализације приказан је на
	Слици~\ref{fig:h2-tradeoff}. Услови са алатима за претрагу налазе се десно и
	ниже, што значи да су скупљи, а да притом не локализују боље.
	
	\begin{figure}[H]
		\centering
		\includegraphics[width=0.95\textwidth]{slike/h2_tradeoff.png}
		\caption{Однос трошка по задатку и стопе локализације по условима}
		\label{fig:h2-tradeoff}
	\end{figure}
	
	\subsection{Провера хипотезе Х2}
	
	Хипотеза Х2 тврди да доступност групе алата за претрагу повећава трошак по
	покретању (Х2а), али истовремено повећава стопу локализације (Х2б). Према
	унапред утврђеном плану, Х2 се сматра потврђеном само ако важе и Х2а и Х2б.
	
	\textbf{Х2а --- потврђена.} Поређен је главни ефекат групе за претрагу,
	услови В и Г наспрам услова А и Б, над упареним разликама по задатку: за
	сваки задатак израчуната је разлика просечног трошка у условима са претрагом
	и просечног трошка у условима без ње. Медијана разлике износи
	$+\$0{,}0064$ (95\% интервал поверења $[+\$0{,}0050,\ +\$0{,}0084]$), а
	Вилкоксонов тест упарених рангова даје $W = 6387$ и
	$p = 5{,}0 \cdot 10^{-27}$, уз величину ефекта $r_{rb} = +0{,}72$. Резултат
	опстаје после Холмове корекције.
	
	Мера трошка рачуна се по ценовнику над свим улазним токенима, без попуста за
	кеширани улаз, па представља горњу границу стварног трошка, што је
	размотрено међу ограничењима рада. Пошто је таква мера детерминистичка
	функција броја токена, поређење услова њоме није изложено утицају
	кеширања.
	
	\textbf{Х2б --- одбачена, и то у супротном смеру.} Разлика стопе
	локализације између услова са претрагом и услова без ње износи $-4{,}67$
	процентних поена, са упареним бутстреп интервалом поверења од 95\%
	$[-8{,}33,\ -1{,}00]$ процентних поена (10\,000 понављања, поновно
	узорковање задатака). Критеријум потврде био је да цео интервал лежи изнад
	нуле, а добијени интервал лежи цео испод нуле. Допунски Мекнемарови
	егзактни тестови дају за пар А--В $b = 31$, $c = 20$ и $p = 0{,}161$, а за
	пар Б--Г $b = 35$, $c = 18$ и $p = 0{,}027$. Ни један не опстаје после
	Холмове корекције примарне фамилије тестова, али оба имају исти смер као и
	интервални резултат.
	
	Овај резултат приказан је и на Слици~\ref{fig:h2-interval}, где се види да
	цео интервал поверења лежи испод нуле.
	
	\begin{figure}[H]
		\centering
		\includegraphics[width=0.95\textwidth]{slike/h2_interval.png}
		\caption{Упарени бутстреп интервал поверења за разлику стопе локализације
			између услова са претрагом и услова без ње}
		\label{fig:h2-interval}
	\end{figure}
	
	Пошто Х2б није потврђена, хипотеза Х2 се одбацује. Штавише, оба њена дела
	показују ефекат супротан очекиваном: додавање алата за претрагу истовремено
	повећава трошак и смањује стопу локализације.
	
	\subsection{Провера хипотезе Х3}
	
	Хипотеза Х3 формулисана је као тврдња о неинфериорности: стопа решених
	задатака у услову Б (\texttt{bash} и \texttt{str\_replace}) није нижа од
	стопе у услову Г (сви алати) за више од пет процентних поена. Маргина је
	утврђена пре оцењивања решености. Критеријум потврде јесте да доња граница
	упареног бутстреп интервала поверења разлике $R_\text{Б} - R_\text{Г}$ лежи
	изнад минус пет процентних поена. Упарени исходи ова два услова приказани су
	у Табели~\ref{tab:h3_parovi}.
	
	\begin{table}[H]
		\caption{Упарени исходи услова Б и Г (300 задатака)}
		\label{tab:h3_parovi}
		\tabelalevo
		\begin{tabular}{|l|l|l|l|}
			\hline
			 & Г решио & Г није решио & Укупно \\
			\hline
			Б решио & 78 & 41 & 119 \\
			\hline
			Б није решио & 6 & 175 & 181 \\
			\hline
			Укупно & 84 & 216 & 300 \\
			\hline
		\end{tabular}
	\end{table}
	
	\textbf{Х3 --- потврђена.} Услов Б решио је 119 од 300 задатака (39,7\%,
	95\% интервал поверења $[34{,}0,\ 45{,}3]$\%), а услов Г 84 задатка
	(28,0\%, интервал $[23{,}0,\ 33{,}3]$\%). Разлика износи $+11{,}67$
	процентних поена, са упареним бутстреп интервалом поверења од 95\%
	$[+7{,}33,\ +16{,}00]$ процентних поена (10\,000 понављања, поновно
	узорковање задатака). Доња граница интервала лежи 12,33 процентна поена
	изнад маргине неинфериорности. Апсолутне стопе решених задатака далеко су од
	нуле, па потврда неинфериорности није последица уског интервала при ниским
	стопама. Стопе решених задатака по условима приказане су на
	Слици~\ref{fig:h3-stope}, интервал поверења разлике на
	Слици~\ref{fig:h3-interval}, а његова бутстреп расподела на
	Слици~\ref{fig:h3-bootstrap}.
	
	\begin{figure}[H]
		\centering
		\includegraphics[width=0.95\textwidth]{slike/h3_stope.png}
		\caption{Стопе решених задатака по условима, са 95\% бутстреп интервалима поверења}
		\label{fig:h3-stope}
	\end{figure}
	
	\begin{figure}[H]
		\centering
		\includegraphics[width=0.95\textwidth]{slike/h3_interval.png}
		\caption{Интервали поверења за разлику стопа решених задатака између услова Б и Г у односу на маргину неинфериорности}
		\label{fig:h3-interval}
	\end{figure}
	
	\begin{figure}[H]
		\centering
		\includegraphics[width=0.95\textwidth]{slike/h3_bootstrap.png}
		\caption{Бутстреп расподела разлике стопа решених задатака између услова Б и Г у примарној анализи}
		\label{fig:h3-bootstrap}
	\end{figure}
	
	Интервал у целини лежи изнад нуле, што указује да услов Б решава више
	задатака од услова Г. Супериорност није била унапред утврђена хипотеза, али
	је у складу са унапред утврђеним секундарним поређењем: егзактни Мекнемаров
	тест за пар Б--Г даје $b = 41$, $c = 6$ и $p = 1{,}8 \cdot 10^{-7}$, и
	опстаје после Холмове корекције фамилије од четири контраста нацрта. Од
	осталих контраста значајан је још само пар В--Г ($b = 32$, $c = 11$,
	$p = 0{,}0019$), док парови А--Б ($p = 0{,}256$) и А--В ($p = 0{,}441$)
	нису.
	
	Анализа осетљивости, ограничена на 256 задатака на којима ни у једном
	услову није дошло до одбијања захтева од стране провајдера, потврђује
	закључак: $R_\text{Б} = 45{,}31\%$, $R_\text{Г} = 31{,}25\%$, разлика
	$+14{,}06$ процентних поена, са интервалом $[+9{,}38,\ +19{,}14]$.
	
	\section{Дискусија резултата}
	
	Одбацивање хипотезе Х1 последица је структуре трошка агента, а не
	неисправности алата. Алат \texttt{str\_replace} ради оно за шта је намењен:
	измена се задаје навођењем сегмента текста уместо поновним исписивањем
	датотеке или састављањем команде, па број излазних токена опада, при чему
	медијана упарене разлике између услова Б и А износи приближно 150 токена по
	задатку. Излазни токени, међутим, чине мање од 10\% трошка. Преостали део
	отпада на улазне токене, јер се у сваком позиву модела шаље цела историја
	разговора, па трошак одређују количина прочитаног кода и број корака, а не
	величина измене, која обично заузима један до два корака на крају
	покретања.
	
	Насупрот малој уштеди стоји фиксни трошак: декларација алата, која по
	мерењу на моделу додаје 223 улазна токена, шаље се у сваком позиву модела,
	што на око 16 корака износи више од 3\,500 улазних токена по покретању,
	неколико пута више од уштеде на излазу. Модел при том алат не користи доследно: у
	102 од 300 покретања услова Б није га позвао ниједном. Упарена разлика
	трошка по покретању између услова Б и А није статистички значајна (медијана
	$+\$0{,}0004$, Вилкоксон $p = 0{,}10$), па алат \texttt{str\_replace} у
	условима без претраге трошак ни не смањује ни значајно не повећава. Главни
	ефекат од $-3{,}1\%$ по покретању потиче из реда нацрта са алатима за
	претрагу, где је истовремено праћен падом стопе решених задатака.
	
	Удруживањем услова прикрива се различито понашање алата у два реда нацрта.
	У условима без алата за претрагу, при прелазу са услова А на Б, трошак по
	решеном задатку готово је непромењен (\$0,0565 наспрам \$0,0585). У
	условима са претрагом, при прелазу са В на Г, покретања постају јефтинија
	(\$0,0352 наспрам \$0,0311 по покретању), али стопа решених пада са 35,0\%
	на 28,0\%, па трошак по решеном задатку расте за приближно 10\% (\$0,1006
	наспрам \$0,1111). Интервал поверења за разлику ова два ефекта обухвата
	нулу ($+\$0{,}0085$, 95\% интервал поверења $[-\$0{,}0099,\ +\$0{,}0302]$),
	па се образац наводи као експлоративан налаз, а не као утврђена интеракција.
	Он је ипак у складу са налазом за хипотезу Х2: највећи скуп алата производи
	највише покретања која исцрпе дозвољени број корака (54) и највише
	изосталих исправки (52).
	
	Трошак по решеном задатку знатно више зависи од алата за претрагу него од
	алата \texttt{str\_replace}: услови без претраге решавају задатак за
	\$0,057--0,059, а услови са претрагом за \$0,101--0,111, дакле приближно
	двоструко скупље. Са становишта цене, избор скупа алата за претрагу важнији
	је од избора алата за измену.
	
	
	Резултат за хипотезу Х2 супротан је полазној претпоставци да наменски алати
	за претрагу помажу агенту да пронађе место грешке. Механизам се назире у
	Табелама~\ref{tab:rezultati} и \ref{tab:ishodi}: у условима са претрагом
	просечан број корака расте са приближно 16 на 22--24, а број покретања која
	исцрпе дозвољени број корака расте са 4--6 на 41--54. Уместо да скрате пут
	до одговарајуће датотеке, ти алати подстичу опширније претраживање, које
	троши буџет корака пре него што агент стигне до измене. Број изосталих исправки
	прати исти образац: 22 у услову А наспрам 52 у услову Г. Услови са
	претрагом троше 1,55 пута више улазних токена, при чему је број токена по
	једном позиву алата готово једнак у сва четири услова, између 6\,058 и
	6\,889. Разлика, дакле, не потиче од обимнијег излаза алата, него од већег
	броја позива алата, који по задатку расте са 15,5 у услову А на 26,3 у
	услову В. Позив алата није исто што и корак: модел у једном кораку може да
	позове више алата, што се у условима са претрагом дешава знатно чешће, а
	последњи корак често нема позив алата, јер модел само завршава одговором.
	
	Овај налаз је у складу са резултатима \textcite{liu2024lostinthemiddle}, по
	којима модели слабије користе информације у средини дугачког контекста, као
	и са тезом \textcite{xia2025agentless} да сложенији систем не доноси нужно
	бољи резултат. Истовремено, он одступа од налаза \textcite{yang2024sweagent},
	где је прилагођени интерфејс, укључујући и претрагу, побољшао успешност у
	односу на рад само са љуском. Разлика може потицати од величине модела,
	пошто је овде коришћен мали модел, коме дужи контекст и већи број могућих
	акција више отежавају него што помажу.
	
	Разлика између услова А и Б није статистички разлучива на мери локализације
	($b = 17$, $c = 15$, $p = 0{,}86$). То није доказ да алат
	\texttt{str\_replace} нема ефекта, него последица тога што мера не може да
	га обухвати: тај алат делује у фази измене датотеке, након што је место
	грешке већ пронађено, док локализација мери искључиво проналажење. Вредност
	тог алата проверава се хипотезом Х1, која се ослања на број решених
	задатака.
	
	Од 300 задатака, на 180 су сва четири услова локализовала место грешке, а на
	35 ни један. Преосталих 85 задатака, на којима се услови разликују, носе сав
	сигнал о утицају скупа алата. Овај однос је важан при тумачењу: удео
	задатака на којима се конфигурације уопште могу разликовати знатно је мањи
	од укупне величине збирке.
	
	Резултат за хипотезу Х3 најјачи је налаз експеримента: најмањи скуп алата
	који садржи наменски алат за измену не само да не заостаје за пуним скупом,
	него решава приближно 42\% више задатака (119 наспрам 84). Од 47 задатака на
	којима се исходи ова два услова разликују, услов Б је успешнији на 41.
	Упарени исходи и исход услова Г на задацима које је решио само услов Б
	приказани су на Слици~\ref{fig:h3-parovi}.
	
	\begin{figure}[H]
		\centering
		\includegraphics[width=0.85\textwidth]{slike/h3_parovi.png}
		\caption{Упарени исходи услова Б и Г и исходи услова Г на задацима које је решио само услов Б}
		\label{fig:h3-parovi}
	\end{figure}
	
	Разлагање стопе решених задатака на узастопне фазе показује где услов Г
	губи. Када произведе исправку, услов Г погађа одговарајућу датотеку једнако
	често као услов Б (86,7\% наспрам 86,2\%), па проналажење места грешке није
	узрок разлике. Разлика настаје у квалитету измене: међу покретањима која
	мењају одговарајућу датотеку, услов Б решава 49,1\% задатака, а услов Г
	38,1\%. Мањи део губитка потиче од изосталих исправки (17,3\% наспрам
	10,3\% покретања), углавном због исцрпљеног броја корака. На 41 задатку који
	је решио само услов Б, агент у услову Г је у 11 случајева исцрпео дозвољени
	број корака без измене, а у 23 случаја изменио одговарајућу датотеку, али
	погрешно. Разлагање по фазама за сва четири услова приказано је на
	Слици~\ref{fig:h3-levak}.
	
	\begin{figure}[H]
		\centering
		\includegraphics[width=0.95\textwidth]{slike/h3_levak.png}
		\caption{Разлагање стопе решених задатака на узастопне фазе по условима}
		\label{fig:h3-levak}
	\end{figure}
	
	Експлоративна анализа транскрипата указује на могући механизам. Агент у
	услову Г најређе покреће код: у просеку 1,85 пута по задатку, наспрам 2,58 у
	услову Б, а пре прве измене, дакле ради репродукције грешке, 0,74 наспрам
	1,19 пута. Код се може покренути само кроз алат \texttt{bash}, а у услову Г
	тај алат чини свега 21\% позива алата, док 54\% отпада на претрагу, па агент
	грешку чешће тумачи читањем кода него што је репродукује. Услов В, који има
	исту групу алата за претрагу, задржава већи удео алата \texttt{bash} (29\%)
	и чешће репродукује грешку (1,29 пута), а на одговарајућој датотеци решава
	45,7\% задатака, приближно као услови А и Б. Пад се, према томе, везује за
	комбинацију оба фактора, што је у складу са секундарним поређењима, у којима
	су значајни само парови Б--Г и В--Г. Бројање покретања кода засновано је на
	хеуристичкој анализи команди и није унапред утврђено, а разлика у стопи
	решених постоји и међу покретањима која су код покренула, па се
	репродукција не може сматрати јединим узроком. Број покретања кода пре и
	после прве измене приказан је на Слици~\ref{fig:h3-reprodukcija}.
	
	\begin{figure}[H]
		\centering
		\includegraphics[width=0.95\textwidth]{slike/h3_reprodukcija.png}
		\caption{Просечан број покретања кода пре и после прве измене по условима}
		\label{fig:h3-reprodukcija}
	\end{figure}
	
	Предност услова Б није ограничена на поједине репозиторијуме: услов Б решио
	је најмање онолико задатака колико и услов Г у свих 12 репозиторијума, а
	више у њих девет, највише у \texttt{django} и \texttt{sympy}, у оба за по
	десет задатака. Поређење по репозиторијумима је експлоративно и без
	корекције за вишеструко тестирање, а исходи по репозиторијумима приказани су
	на Слици~\ref{fig:h3-repo}. Пошто су покретања прекинута одбијањем
	захтева знатно чешћа у услову Б (29 наспрам три), а у примарној анализи се
	рачунају као нерешена, примарни резултат је у том погледу конзервативан:
	искључивањем тих задатака разлика се повећава.
	
	\begin{figure}[H]
		\centering
		\includegraphics[width=0.85\textwidth]{slike/h3_repo.png}
		\caption{Број решених задатака у условима Б и Г по репозиторијумима}
		\label{fig:h3-repo}
	\end{figure}
	
	Посматрани заједно, резултати делимично подржавају полазну тезу рада. Од
	три хипотезе потврђена је само Х3, али је она најзначајнија за практичну
	препоруку: комбинација алата \texttt{bash} и \texttt{str\_replace} решава
	највише задатака од свих услова (39,7\%), уз трошак по решеном задатку који
	се статистички не разликује од најјефтинијег услова А (\$0,0585 наспрам
	\$0,0565). Други део тезе није потврђен, јер пун скуп алата не доноси
	предност ни при лоцирању грешке, него је и по трошку и по стопи
	локализације лошији од мањих скупова.
	
	\section{Ограничења}
	
	Резултате рада треба тумачити у светлу следећих ограничења:
	\begin{itemize}
		\item \textbf{Један модел.} Експеримент је спроведен само са моделом
		\texttt{gpt-5.4-nano}, па се резултати не могу непосредно уопштити на
		јаче моделе. Смер тог утицаја не може се унапред претпоставити:
		очекивање да слабији модели имају више користи од помоћних алата није
		потврђено, јер је пун скуп алата малом моделу одмогао.
		Посебно, није проверено да ли се разлика између услова Б и Г одржава
		код модела веће способности, који би алате за претрагу могли користити
		делотворније.
		\item \textbf{Природа скупа SWE-bench Lite.} Скуп садржи искључиво
		самосталне исправке у једној датотеци са највише три измењена блока, што
		је тип измене у коме \texttt{str\_replace} има највећу предност. Код
		измена кроз више датотека ефекат може бити мањи.
		\item \textbf{Квалитет скупа података.} Део задатака садржи решење у
		опису проблема, а део тестова није довољан да открије нетачне исправке
		\parencite{aleithan2024swebenchplus}. Постоји и ризик да су задаци били
		део података за тренирање модела. Резултати зато важе у оквиру
		бенчмарка.
		\item \textbf{Заштитни механизми провајдера.} Модел је коришћен преко
		Azure OpenAI сервиса, чији заштитни механизми (енг. \textit{guardrails})
		проверавају садржај сваког захтева. Пошто су захтеви садржали изворни
		код и излазе команди, класификатор је део њих означавао као неприхватљив
		садржај, иако то нису били, и враћао грешку \texttt{invalid\_prompt}.
		Понављање захтева отклонило је део ових случајева, али не све, па
		одређени број задатака није могао да буде решен из разлога који не
		зависе ни од модела ни од скупа алата. Стопа таквих грешака извештава се
		по услову.
		\item \textbf{Недостајући подаци нису независни од услова.} Одбијања
		описана у претходној ставци нису равномерно распоређена: у условима без
		алата за претрагу било их је 52 (23 у услову А и 29 у услову Б), а у
		условима са претрагом свега девет (шест у услову В и три у услову Г).
		Разлог је то што агент без наменских алата више необрађеног кода
		пропушта кроз команде љуске. Недостајући подаци тиме корелирају са
		независном променљивом, па се њихова стопа извештава по услову, а
		налази се проверавају анализом осетљивости на задацима који су завршени
		у сва четири услова. У примарној анализи ова покретања улазе као исход
		услова: трошак који су до тренутка одбијања произвела урачунава се, а
		задатак није решен, што оптерећује трошак по решеном задатку у условима
		без претраге. Анализа осетљивости даје исти закључак за хипотезе Х1 и Х3.
		\item \textbf{Једно покретање по услову.} Није могуће проценити
		варијабилност резултата истог задатка у истом услову, а излаз модела
		није потпуно детерминистички ни при температури 0.
		\item \textbf{Агент ради „наслепо“.} Због јединствене слике окружења,
		описане у поглављу~4, агент не може да покрене тестове пројекта нити да
		репродукује грешку, него се ослања само на читање кода и опис проблема.
		То је значајно одступање од упоредивих система, попут SWE-agent
		\parencite{yang2024sweagent} и Agentless \parencite{xia2025agentless},
		који агенту омогућавају покретање тестова и проверу исправке, па је
		стопа решених задатака вероватно нижа него што би иначе била. Одлука
		важи подједнако у свим условима, тако да не утиче на поређење услова.
		\item \textbf{Усклађеност алата са обуком модела.} Позивање алата је
		научена способност \parencite{schick2023toolformer}, па оно зависи и од
		тога колико су назив и шема понуђеног алата слични онима које је модел
		видео током обуке. Алати развијени за овај рад не морају се поклапати са
		тим обрасцима, на шта указује и то што је модел у појединим покретањима
		позивао алате који у овом систему не постоје. Утицај овог чиниоца није
		мерљив овим нацртом, а у складу је са налазом
		\textcite{yang2024sweagent} да дизајн интерфејса битно утиче на
		успешност агента.
		\item \textbf{Резултат важи за ову имплементацију алата.} Измерени ефекат
		приписује се скупу алата, али зависи и од тога како је сваки алат
		написан. Мерења говоре против објашњења да су алати неисправни:
		\texttt{search} и \texttt{find\_file} нису вратили ниједну грешку у
		8\,467 позива, \texttt{read\_file} у 0,6\% позива, а број токена по
		једном позиву алата готово је једнак у сва четири услова (између 6\,058 и
		6\,889), што искључује и претпоставку да алати за претрагу претрпавају
		контекст. Ипак, поједини детаљи њиховог нацрта могли су допринети већем
		броју позива алата: \texttt{search} враћа највише 50 погодака без могућности
		страничења, без редова контекста око поготка и без ознаке да је излаз
		одсечен, због чега модел не разликује потпун од непотпуног резултата и
		мора да допуни увид додатним позивима. Другачији нацрт исте
		функционалности могао би дати другачији исход, па закључак треба читати
		као тврдњу о овој површини алата, а не о алатима за претрагу уопште.
		\item \textbf{Неједнако обавезно образложење потеза.} Декларација алата
		\texttt{bash} захтева додатно поље у коме модел пре сваког позива
		образлаже поступак, док остали алати такво поље немају. Пошто се
		\texttt{bash} користи неједнако међу условима (100\% позива алата у услову А,
		21\% у услову Г), и то образлагање пада неједнако: приближно 356 токена
		по задатку у услову А наспрам 98 у услову Г. Реч је о подстицају на
		разложено закључивање примењеном неравномерно по условима, чији смер није
		једнозначан, јер може подједнако помагати размишљању и трошити токене.
		Удео у укупној потрошњи је мали (испод 0,5\%), а утицај на стопу
		локализације и решених задатака није мерљив овим нацртом.
		\item \textbf{Мера локализације није мера исправности.} Локализација
		утврђује само да исправка агента дира бар једну датотеку коју дира и
		референтна исправка, не и да је измена исправна. Закључак о хипотези Х2
		стога важи за фазу проналажења места грешке, док се исправност проверава
		хипотезама Х1 и Х3, на основу резултата званичног евалуатора.
		\item \textbf{Мера трошка не обухвата кеширање.} Трошак се рачуна по
		ценовнику над свим улазним токенима, без попуста за кеширани улаз, па
		представља горњу границу стварног трошка; број кешираних токена није
		бележен током извођења, те се попуст не може накнадно применити. Улазни
		токени чине 94\% обрачуна, а стварни трошак лежи између \$5,19 (уз
		потпуно кеширање) и \$33,13 (без кеширања), при чему се стварни рачун
		провајдера налазио близу доње границе. Пошто је употребљена мера
		детерминистичка функција броја токена, поређење услова њоме није
		изложено утицају кеширања. Ипак, услови са претрагом имају дуже путање,
		па и вероватно виши удео кеширања, услед чега би разлика у стварном
		трошку могла бити мања од пријављене. Мера је изабрана зато што је
		поновљива и независна од провајдера, за разлику од стварне стопе
		кеширања, која зависи од распореда и времена извођења.
		\item \textbf{Трошак по решеном задатку је однос.} Пошто трошак по
		решеном задатку не постоји за појединачни нерешен задатак, мера се не
		може поставити у упарени тест по задатку. Интервал поверења добијен је
		бутстреп методом, а допунски Вилкоксонов тест односи се на трошак по
		покретању, који није иста мера. Однос збирова осетљив је на мали број
		скупих покретања, што је при 203 и 216 решених задатака по групи
		ублажено, али не и отклоњено.
		\item \textbf{Маргина неинфериорности задата је одлуком аутора.} Вредност
		од пет процентних поена није изведена из спољашњег стандарда, већ
		представља процену губитка који би био практично занемарљив у односу на
		смањење сложености скупа алата. Маргина је утврђена пре оцењивања
		решености. Добијени резултат није осетљив на њен избор, јер доња граница
		интервала лежи изнад нуле, тако да би закључак важио за сваку ненегативну
		маргину.
		\item \textbf{Несигурни исходи оцењивања.} За мали број задатака званични
		евалуатор није произвео резултат тестова, па се задатак рачуна као
		нерешен. У условима Б и Г таквих је ниједан односно два задатка, а
		њихово прерачунавање у решене у услову Г смањило би разлику за 0,67
		процентних поена, без утицаја на закључак.
		\item \textbf{Тренутак утврђивања хипотеза.} Хипотезе и план статистичке
		анализе уписани су током извођења експеримента: услов А је био завршен у
		целости, услов Б делимично, а услови В и Г нису били покренути.
		Утврђивање је, међутим, извршено пре било каквог оцењивања решености и
		пре првог израчунавања локализације, тако да ни један исход који улази у
		проверу хипотеза није био познат у тренутку уписа.
	\end{itemize}
	
	% ============================= 6. ЗАКЉУЧАК ==================================
	\chapter{Закључак}
	
	Предмет овог рада било је питање у којој мери скуп алата који агент за
	програмирање има на располагању утиче на његову успешност и трошак,
	независно од избора језичког модела. Питање је испитано на практичној
	имплементацији агента и на целом \texttt{test} подскупу збирке SWE-bench
	Lite, у нацрту $2 \times 2$ који раздваја два фактора: наменски алат за
	измену датотека (\texttt{str\_replace}) и групу алата за претрагу и читање
	кода (\texttt{search}, \texttt{read\_file} и \texttt{find\_file}). Сва четири
	услова покренута су над истих 300 задатака, са истим моделом и истим
	системским упутством, а хипотезе, маргина неинфериорности и план
	статистичке анализе утврђени су пре оцењивања решености.
	
	\section{Главни налази}
	
	Резултати не потврђују интуицију да агент са богатијим скупом алата ради
	боље. Најмањи скуп алата који садржи наменски алат за измену, \texttt{bash}
	и \texttt{str\_replace}, решио је 119 од 300 задатака (39,7\%), а пун скуп
	алата 84 (28,0\%). Разлика од 11,7 процентних поена, са 95\% интервалом
	поверења $[+7{,}3,\ +16{,}0]$, потврђује хипотезу Х3 о неинфериорности мањег
	скупа и указује да мањи скуп решава више задатака. Мањи скуп је истовремено
	и јефтинији: трошак по решеном задатку износи \$0,059 наспрам \$0,111.
	
	Хипотеза Х2, према којој алати за претрагу повећавају трошак, али побољшавају
	проналажење места грешке, потврђена је само у првом делу. Услови са
	претрагом троше знатно више ($p = 5{,}0 \cdot 10^{-27}$, приближно 1,55 пута
	више улазних токена), али стопу локализације не повећавају, него је смањују
	за 4,7 процентних поена, са 95\% интервалом поверења $[-8{,}3,\ -1{,}0]$.
	Додатни трошак не потиче од обимнијег излаза алата, него од већег броја
	позива алата и корака: број корака расте са приближно 16 на 22--24, а број
	покретања која исцрпе дозвољени буџет са 4--6 на 41--54.
	
	Хипотеза Х1, према којој \texttt{str\_replace} смањује трошак по решеном
	задатку, одбачена је: разлика износи $+3{,}1\%$, уз интервал поверења који
	обухвата нулу. Алат смањује број излазних токена, али излазни токени чине
	мање од 10\% трошка. Трошак агента одређује количина прочитаног контекста,
	који се шаље у сваком позиву модела, а не величина измене, која заузима један
	до два корака на крају рада. Уштеда на излазу мања је од цене декларације
	алата, која се такође шаље у сваком позиву модела.
	
	\section{Тумачење}
	
	Три резултата указују на исти образац: за мали модел свака додатна
	могућност има цену која није садржана у самом алату. Алати за претрагу нису
	произвели неисправан или преобиман излаз, јер је број токена по позиву
	алата готово једнак у свим условима, али су променили начин рада агента. У
	пуном скупу алата претрага чини више од половине позива алата, а
	\texttt{bash} свега петину, и агент најређе покреће код, нарочито пре прве
	измене (0,74 наспрам 1,19 пута по задатку у услову Б). Место грешке при
	томе проналази једнако успешно, али га ређе исправно мења: међу покретањима
	која мењају одговарајућу датотеку, пун скуп алата решава 38,1\% задатака, а
	услов Б 49,1\%. Могуће објашњење јесте да агент грешку чешће тумачи читањем
	кода него што је репродукује. Та веза утврђена је експлоративно и не може се
	сматрати доказаном, али је у складу са чињеницом да ниједан од два фактора
	засебно не доводи до значајног пада, већ тек њихова комбинација.
	
	Шире посматрано, резултати указују да се избор алата не може свести на
	питање да ли алат ради оно што му је намена. И \texttt{str\_replace} и алати
	за претрагу раде исправно, а њихов утицај на агента одређују структура
	трошка и промена понашања коју изазивају. Те две ствари не виде се из описа
	алата, већ тек из мерења.
	
	\section{Практичне препоруке}
	
	За агента заснованог на малом моделу, резултати упућују на следеће:
	\begin{itemize}
		\item почети од минималног скупа алата (\texttt{bash} и алат за измену) и
		сваки додатни алат уводити тек када мерење покаже корист;
		\item трошак процењивати по решеном задатку, а не по покретању, и уз њега
		пратити број корака и удео покретања која исцрпе буџет;
		\item оптимизацију трошка усмерити на количину контекста који агент чита,
		а не на излаз модела, јер улазни токени чине више од 90\% трошка;
		\item при увођењу алата за претрагу пратити да ли агент и даље проверава
		своје претпоставке покретањем кода.
	\end{itemize}
	
	\section{Ограничења}
	
	Закључци важе у границама нацрта. Сви услови покренути су са једним, малим
	моделом, при температури 0 и са једним покретањем по задатку, па није
	проверено да ли се образац одржава код модела веће способности, који би
	алате за претрагу могли користити другачије; смер тог утицаја не може се
	унапред претпоставити. Измерени ефекат односи се на ову имплементацију
	алата, а не на алате за претрагу уопште. Трошак је рачунат без попуста за
	кеширани улаз, па представља горњу границу стварног трошка. Део покретања
	прекинут је одбијањем захтева од стране провајдера, неравномерно по
	условима, а анализа осетљивости која те задатке искључује даје исте
	закључке за хипотезе Х1 и Х3. Механизам описан у тумачењу утврђен је
	експлоративном анализом транскрипата и није био део унапред утврђеног плана.
	
	\section{Даљи рад}
	
	Најнепосреднији наставак јесте понављање експеримента са моделом веће
	способности, макар на подскупу задатака, чиме би се проверило да ли је налаз
	својство малог модела или скупа алата. Поновљена покретања по задатку
	омогућила би да се процени колики део неслагања између услова потиче од
	случајности. Посебно је занимљиво испитати да ли се губитак пуног скупа
	алата може отклонити променом нацрта самих алата, на пример претрагом са
	страничењем и контекстом око погодака, или изричитим подстицањем агента да
	грешку репродукује пре измене. Коначно, мерење стварне стопе кеширања и
	укључивање алата за покретање тестова уз приступ мрежи проширили би слику о
	томе како скуп алата утиче на однос цене и успешности.
	
	% -----------------------------------------------------------------------------
	% УПОТРЕБА АЛАТА ЗАСНОВАНИХ НА ВЕШТАЧКОЈ ИНТЕЛИГЕНЦИЈИ
	% -----------------------------------------------------------------------------
	%\newpage
	%\phantomsection
	%\addcontentsline{toc}{chapter}{Употреба алата заснованих на вештачкој интелигенцији}
	%\chapter*{Употреба алата заснованих на вештачкој интелигенцији}
	%
	%У изради овог рада коришћени су алати засновани на великим језичким
	%моделима, и то на два начина.
	%
	%Прво, велики језички модел представља предмет истраживања: агентски систем
	%развијен за потребе рада користи модел \texttt{gpt-5.4-nano}, а сви
	%резултати приказани у раду добијени су покретањем тог система, како је
	%описано у поглављима~4 и~5.
	%
	%Друго, при писању рада коришћен је асистент заснован на великом језичком
	%моделу, и то за обликовање и језичко дотеривање текста, проналажење и
	%проверу библиографских података, као и за помоћ при обликовању документа у
	%систему \LaTeX. Идеја рада, имплементација агентског система, поставка и
	%извођење експеримента, обрада резултата и закључци представљају самосталан
	%рад аутора. Аутор је проверио све податке, тврдње и референце наведене у
	%раду и преузима одговорност за њихову тачност.
	
	% -----------------------------------------------------------------------------
	% ЛИТЕРАТУРА
	% -----------------------------------------------------------------------------
	\printbibliography[title={Литература}, heading=bibintoc]


\end{document}